\begin{proof}[Proof of \cref{obj:prop:bounded-submodel-comparison}]
Fix the admissible tuple \(v=(p,\alpha,\beta,\gamma)\) and its matched smoothness tuple \(w=(\alpha,\beta,\gamma)\).

\begin{enumerate}
\item We first establish the integrated conversion identity. Write
\[
\mathcal O=[0,1]\times\{0,1\}\times\mathbb R
\]
for the Borel record space. For a law \(P\) in either
\(\mathcal M_v\) or \(\mathcal M_{(2,\alpha,\beta,\gamma)}\), let
\(Q_{a_{\mathrm{arm}},P}(\,\cdot\mid x)\) be its conditional outcome probability kernel in arm \(a_{\mathrm{arm}}\in\{0,1\}\). Thus
\[
\begin{split}
P(dx,\{a_{\mathrm{arm}}\},dy)
&=\lambda(dx)e_P(x)^{a_{\mathrm{arm}}}
  (1-e_P(x))^{1-a_{\mathrm{arm}}}
  Q_{a_{\mathrm{arm}},P}(dy\mid x),\\
m_{a_{\mathrm{arm}},P}(x)
&=m_{0,P}(x)+a_{\mathrm{arm}}\tau_P(x)
 =\int y\,Q_{a_{\mathrm{arm}},P}(dy\mid x)
 \qquad\text{for \(\lambda\)-almost every \(x\)}.
\end{split}
\]
The mean bounds give
\[
|m_{0,P}(x)|\le\frac12,\qquad
|m_{1,P}(x)|
\le |m_{0,P}(x)|+|\tau_P(x)|\le1.
\]
Define the signed-binary conversion \(\mathcal C_{\mathrm{bin}}(P)\) by
\[
\begin{split}
Q^{\mathrm{bin}}_{a_{\mathrm{arm}},P}(dy\mid x)
&=\frac{1+m_{a_{\mathrm{arm}},P}(x)}2\,\delta_1(dy)
 +\frac{1-m_{a_{\mathrm{arm}},P}(x)}2\,\delta_{-1}(dy),\\
\mathcal C_{\mathrm{bin}}(P)(dx,\{a_{\mathrm{arm}}\},dy)
&=\lambda(dx)e_P(x)^{a_{\mathrm{arm}}}
 (1-e_P(x))^{1-a_{\mathrm{arm}}}
 Q^{\mathrm{bin}}_{a_{\mathrm{arm}},P}(dy\mid x).
\end{split}
\]
The two arm weights are nonnegative and sum to one, and
\[
\int y\,Q^{\mathrm{bin}}_{a_{\mathrm{arm}},P}(dy\mid x)
=m_{a_{\mathrm{arm}},P}(x),\qquad
\int |y|^2\,Q^{\mathrm{bin}}_{a_{\mathrm{arm}},P}(dy\mid x)=1.
\]
Consequently the conversion has signed-binary support and retains the uniform design, overlap, smoothness restrictions, and mean bounds. In particular,
\[
\begin{gathered}
\mathcal C_{\mathrm{bin}}(P)\in\mathcal M_w^{\mathrm{bin}},\qquad
\bigl(\mathcal C_{\mathrm{bin}}(P)\bigr)_{X,A}=P_{X,A},\\
e_{\mathcal C_{\mathrm{bin}}(P)}=e_P,\qquad
m_{0,\mathcal C_{\mathrm{bin}}(P)}=m_{0,P},\qquad
\tau_{\mathcal C_{\mathrm{bin}}(P)}=\tau_P.
\end{gathered}
\]
Substitution into the definitions of distance and constancy yields
\[
d(\mathcal C_{\mathrm{bin}}(P))=d(P),
\]
and
\[
\begin{split}
&\exists c\in[-1/2,1/2]\ 
  \forall x,\ \tau_{\mathcal C_{\mathrm{bin}}(P)}(x)=c\\
&\hspace{35mm}\Longleftrightarrow
\exists c\in[-1/2,1/2]\ 
  \forall x,\ \tau_P(x)=c.
\end{split}
\]
% lean: CausalSmith.Stat.FinitepHomogeneityDensegamma.binaryConversion_primitives
% lean: CausalSmith.Stat.FinitepHomogeneityDensegamma.binaryConversion_inModel
% lean: CausalSmith.Stat.FinitepHomogeneityDensegamma.binaryConversion_preserves_effect

To define record randomization on the whole record space, set
\[
\ell_1(y)=\max\{-1,\min\{1,y\}\}\qquad(y\in\mathbb R)
\]
and define the probability kernel
\[
\mathsf K_{\mathrm{sign}}((x,a_{\mathrm{obs}},y),\cdot)
=
\frac{1+\ell_1(y)}2\,\delta_{(x,a_{\mathrm{obs}},1)}
+\frac{1-\ell_1(y)}2\,\delta_{(x,a_{\mathrm{obs}},-1)},
\qquad (x,a_{\mathrm{obs}},y)\in\mathcal O.
\]
For \(P\in\mathcal M_w^{\mathrm b}\), uniform design and fixed overlap transfer the observed outcome bound to both conditional arm kernels:
\[
Q_{a_{\mathrm{arm}},P}([-1,1]\mid x)=1
\qquad
\text{for \(\lambda\)-almost every \(x\), for both arms}.
\]
Indeed, disintegrating \(P(|Y|>1)=0\) gives
\[
\int
\left[
(1-e_P(x))Q_{0,P}(\{|y|>1\}\mid x)
+e_P(x)Q_{1,P}(\{|y|>1\}\mid x)
\right]\lambda(dx)=0.
\]
Both summands are nonnegative, and both arm weights are at least \(1/4\). Each conditional probability therefore vanishes almost everywhere. Since \(\ell_1(y)=y\) on \([-1,1]\), for \(b_{\mathrm{sign}}\in\{-1,1\}\),
\[
\int
\frac{1+b_{\mathrm{sign}}\ell_1(y)}2\,
Q_{a_{\mathrm{arm}},P}(dy\mid x)
=
\frac{1+b_{\mathrm{sign}}m_{a_{\mathrm{arm}},P}(x)}2
\qquad\text{for \(\lambda\)-almost every \(x\)}.
\]
Integrating the two atomic weights and then the treatment mixture proves the one-record identity
\[
\int_{\mathcal O}\mathsf K_{\mathrm{sign}}(o,\cdot)\,P(do)
=\mathcal C_{\mathrm{bin}}(P).
\]
% lean: CausalSmith.Stat.FinitepHomogeneityDensegamma.binaryRecordKernel_comp

Now let \(n\in\mathbb N\), let \(\phi\) be a measurable \([0,1]\)-valued test, and write its input as
\[
\boldsymbol o=(o_i)_{i=1}^n\in\mathcal O^n,\qquad
o_i=(x_i,a_{\mathrm{obs},i},y_i),\qquad
u_{\mathrm{seed}}\in[0,1].
\]
Define its integrated rejection map by
\[
\widetilde\phi(\boldsymbol o,u_{\mathrm{seed}})
=
\sum_{\boldsymbol b_{\mathrm{sign}}\in\{-1,1\}^n}
\phi\!\left(
 ((x_i,a_{\mathrm{obs},i},b_{\mathrm{sign},i}))_{i=1}^n,
 u_{\mathrm{seed}}
\right)
\prod_{i=1}^n
\frac{1+b_{\mathrm{sign},i}\ell_1(y_i)}2,
\qquad
\boldsymbol b_{\mathrm{sign}}
=(b_{\mathrm{sign},i})_{i=1}^n.
\]
Its weights are nonnegative and satisfy
\[
\sum_{\boldsymbol b_{\mathrm{sign}}\in\{-1,1\}^n}
\prod_{i=1}^n
\frac{1+b_{\mathrm{sign},i}\ell_1(y_i)}2
=
\prod_{i=1}^n
\left(
\frac{1+\ell_1(y_i)}2+\frac{1-\ell_1(y_i)}2
\right)=1.
\]
Thus \(\widetilde\phi\) is measurable and takes values in \([0,1]\). For \(n=0\), the sign-vector set contains the single empty vector and the empty product is one, so the same definition applies.

The finite sum is the integral of \(\phi\) against the independent product of the record kernels. Independence and the one-record identity imply
\[
\int_{\mathcal O^n}
\bigotimes_{i=1}^n\mathsf K_{\mathrm{sign}}(o_i,\cdot)\,
P^{\otimes n}(d\boldsymbol o)
=
\mathcal C_{\mathrm{bin}}(P)^{\otimes n}.
\]
This follows by factorization on measurable rectangles; the empty product gives the identity for \(n=0\). Hence, for every fixed \(u_{\mathrm{seed}}\),
\[
\int_{\mathcal O^n}
\widetilde\phi(\boldsymbol o,u_{\mathrm{seed}})
\,P^{\otimes n}(d\boldsymbol o)
=
\int_{\mathcal O^n}
\phi(\boldsymbol o,u_{\mathrm{seed}})
\,\mathcal C_{\mathrm{bin}}(P)^{\otimes n}(d\boldsymbol o).
\]
Integrating in the public seed and using the experiment of
\cref{obj:def:original-record-experiment} gives
\[
\mathsf R_n(P,\widetilde\phi)
=
\mathsf R_n(\mathcal C_{\mathrm{bin}}(P),\phi).
\]
Boundedness of the tests justifies the integral exchanges. On bounded records, the kernel uses exactly the probabilities \((1+Y)/2\) and \((1-Y)/2\) stated in the proposition.
% lean: CausalSmith.Stat.FinitepHomogeneityDensegamma.binaryRandomization_rejectProb

\item We next compare the binary and bounded testing risks. In the infimum--supremum comparisons, a supremum of type-II errors over an empty alternative set is assigned the value zero. The test
\[
\phi_{\mathrm{zero}}(\boldsymbol o,u_{\mathrm{seed}})=0
\]
is level-valid for every null class. The sets of level-valid tests are therefore nonempty, and
\[
0\le 1-\mathsf R_n(P,\phi)\le1
\]
bounds every individual type-II error and every corresponding worst-case error.

Signed-binary support implies bounded support. Keeping the remaining restrictions and the constancy condition gives
\[
\mathcal M_w^{\mathrm{bin}}\subseteq\mathcal M_w^{\mathrm b},
\qquad
H_0^{\mathrm{bin}}(w)\subseteq H_0^{\mathrm b}(w).
\]
Every bounded-null level-valid test is consequently binary-null level-valid, with
\[
\sup_{\substack{P\in\mathcal M_w^{\mathrm{bin}}\\d(P)\ge r}}
\{1-\mathsf R_n(P,\phi)\}
\le
\sup_{\substack{P\in\mathcal M_w^{\mathrm b}\\d(P)\ge r}}
\{1-\mathsf R_n(P,\phi)\}.
\]
For a nonempty binary alternative set, this is the supremum comparison under inclusion; for an empty binary alternative set, it follows from the zero convention and nonnegativity. Taking infima gives
\[
B_n^{\mathrm{bin}}(w,r)\le B_n^{\mathrm b}(w,r).
\]
% lean: CausalSmith.Stat.FinitepHomogeneityDensegamma.binaryTestingRisk_le_boundedTestingRisk

Conversely, take any binary-null level-valid test \(\phi\). For every \(P\in H_0^{\mathrm b}(w)\), the conversion belongs to \(H_0^{\mathrm{bin}}(w)\), so the preceding identity gives
\[
\mathsf R_n(P,\widetilde\phi)
=
\mathsf R_n(\mathcal C_{\mathrm{bin}}(P),\phi)
\le\frac1{10}.
\]
For a bounded alternative \(P\) with \(d(P)\ge r\), its conversion is a binary alternative with the same distance, and
\[
1-\mathsf R_n(P,\widetilde\phi)
=
1-\mathsf R_n(\mathcal C_{\mathrm{bin}}(P),\phi)
\le
\sup_{\substack{P\in\mathcal M_w^{\mathrm{bin}}\\d(P)\ge r}}
\{1-\mathsf R_n(P,\phi)\}.
\]
If the bounded alternative set is nonempty, taking its supremum preserves this inequality; the binary supremum is bounded above by one. If the bounded alternative set is empty, its worst-case error is zero and the binary worst-case error is nonnegative. Thus in either case,
\[
\sup_{\substack{P\in\mathcal M_w^{\mathrm b}\\d(P)\ge r}}
\{1-\mathsf R_n(P,\widetilde\phi)\}
\le
\sup_{\substack{P\in\mathcal M_w^{\mathrm{bin}}\\d(P)\ge r}}
\{1-\mathsf R_n(P,\phi)\}.
\]
The bounded-model infimum is at most the left-hand side for each such \(\phi\). Taking the binary-model infimum proves
\[
B_n^{\mathrm b}(w,r)\le B_n^{\mathrm{bin}}(w,r),
\qquad
B_n^{\mathrm{bin}}(w,r)=B_n^{\mathrm b}(w,r).
\]
% lean: CausalSmith.Stat.FinitepHomogeneityDensegamma.boundedTestingRisk_eq_binary_of_conversion

\item For the remaining model inclusion, take \(P\in\mathcal M_w^{\mathrm b}\). The conditional support established above and the nonnegative moment exponent supplied by admissibility give
\[
\int |y|^p\,Q_{a_{\mathrm{arm}},P}(dy\mid x)
\le
\int 1\,Q_{a_{\mathrm{arm}},P}(dy\mid x)
=1\le10
\qquad\text{for \(\lambda\)-almost every \(x\)}.
\]
All the other primitive restrictions are identical at the matched smoothness tuple. Therefore
\[
\mathcal M_w^{\mathrm{bin}}
\subseteq\mathcal M_w^{\mathrm b}
\subseteq\mathcal M_v,
\qquad
H_0^{\mathrm b}(w)\subseteq H_0(v).
\]
% lean: CausalSmith.Stat.FinitepHomogeneityDensegamma.boundedModel_inModel

\item The distance suprema are equal because the attained distance sets are equal. The binary-to-bounded inclusion gives one direction, and converting each bounded law gives the other:
\[
\{d(P):P\in\mathcal M_w^{\mathrm{bin}}\}
=
\{d(P):P\in\mathcal M_w^{\mathrm b}\}.
\]
For comparison with the full model, convert each bounded law. Its converted arm kernels also satisfy
\[
\int |y|^p\,Q^{\mathrm{bin}}_{a_{\mathrm{arm}},P}(dy\mid x)
=1\le10,
\]
so the conversion belongs to \(\mathcal M_v\) and has the same distance. Conversely, the conversion of every \(P\in\mathcal M_v\) belongs to the binary model, hence to the bounded model, and retains \(d(P)\). Thus
\[
\{d(P):P\in\mathcal M_w^{\mathrm b}\}
=
\{d(P):P\in\mathcal M_v\}.
\]
Taking suprema yields
\[
D_w^{\mathrm{bin}}=D_w^{\mathrm b}=D_v.
\]
Finally, \cref{obj:lem:causal-nonempty} supplies the boundedness of the full-model distance set and the witness lower bound
\[
d_0\le D_v\le\frac12.
\]
% lean: CausalSmith.Stat.FinitepHomogeneityDensegamma.maxDistBinary_eq_maxDistBounded
% lean: CausalSmith.Stat.FinitepHomogeneityDensegamma.maxDistBounded_eq_maxDist
% lean: CausalSmith.Stat.FinitepHomogeneityDensegamma.model_distance_lower

\item A test level-valid for \(H_0(v)\) is level-valid for \(H_0^{\mathrm b}(w)\), and model inclusion gives
\[
\sup_{\substack{P\in\mathcal M_w^{\mathrm b}\\d(P)\ge r}}
\{1-\mathsf R_n(P,\phi)\}
\le
\sup_{\substack{P\in\mathcal M_v\\d(P)\ge r}}
\{1-\mathsf R_n(P,\phi)\}.
\]
For a nonempty bounded alternative set, this follows by inclusion; for an empty set, it follows from the zero convention and nonnegativity of the full-model error. Taking infima over full-null level-valid tests, and allowing all bounded-null level-valid tests on the left, proves
\[
B_n^{\mathrm b}(w,r)\le B_n(v,r).
\]
Together with the equality already established, this gives, for \(n\ge2\) and \(0<r<D_v\),
\[
B_n^{\mathrm{bin}}(w,r)=B_n^{\mathrm b}(w,r)\le B_n(v,r).
\]
% lean: CausalSmith.Stat.FinitepHomogeneityDensegamma.boundedTestingRisk_le_testingRisk

\item We establish strict positivity of the binary capped radius by a small cosine alternative. For \(t_{\mathrm{cos}}\in[0,1]\), define the deterministic law and its conversion by
\[
\begin{gathered}
P^{\mathrm{det}}_{\mathrm{cos},t_{\mathrm{cos}}}:
\quad X\sim\lambda,\qquad
\Pr(A=1\mid X=x)=\frac12,\qquad
Y=A\,\frac{t_{\mathrm{cos}}}{8}\cos(2\pi X),\\
P_{\mathrm{cos},t_{\mathrm{cos}}}
=\mathcal C_{\mathrm{bin}}
 \bigl(P^{\mathrm{det}}_{\mathrm{cos},t_{\mathrm{cos}}}\bigr).
\end{gathered}
\]
Their primitive functions are
\[
e_{P_{\mathrm{cos},t_{\mathrm{cos}}}}(x)=\frac12,\qquad
m_{0,P_{\mathrm{cos},t_{\mathrm{cos}}}}(x)=0,\qquad
\tau_{P_{\mathrm{cos},t_{\mathrm{cos}}}}(x)
=\frac{t_{\mathrm{cos}}}{8}\cos(2\pi x).
\]
The propensity and baseline are constant. For the effect, the cosine increment bound and \(|x-z|\le1\) imply
\[
\begin{split}
|\tau_{P_{\mathrm{cos},t_{\mathrm{cos}}}}(x)|
&\le\frac{t_{\mathrm{cos}}}{8}\le\frac12,\\
|\tau_{P_{\mathrm{cos},t_{\mathrm{cos}}}}(x)
 -\tau_{P_{\mathrm{cos},t_{\mathrm{cos}}}}(z)|
&\le t_{\mathrm{cos}}\frac{2\pi}{8}|x-z|
 \le t_{\mathrm{cos}}\,20|x-z|^\gamma
 \le20|x-z|^\gamma.
\end{split}
\]
The deterministic outcomes have absolute value at most \(1/8\), so their conditional second moments are at most one. These bounds verify membership of the deterministic law in
\(\mathcal M_{(2,\alpha,\beta,\gamma)}\); conversion then gives
\[
P_{\mathrm{cos},t_{\mathrm{cos}}}\in\mathcal M_w^{\mathrm{bin}},
\qquad
P_{\mathrm{cos},0}\in H_0^{\mathrm{bin}}(w).
\]
Moreover,
\[
\int\tau_{P_{\mathrm{cos},t_{\mathrm{cos}}}}\,d\lambda=0,\qquad
\int\tau_{P_{\mathrm{cos},t_{\mathrm{cos}}}}^2\,d\lambda
=\frac{t_{\mathrm{cos}}^2}{128},
\qquad
d(P_{\mathrm{cos},t_{\mathrm{cos}}})=t_{\mathrm{cos}}d_0.
\]
% lean: CausalSmith.Stat.FinitepHomogeneityDensegamma.smallCosineLaw_inBinaryModel
% lean: CausalSmith.Stat.FinitepHomogeneityDensegamma.smallCosineLaw_distance

For probability measures \(\mu_{\mathrm{tv}},\nu_{\mathrm{tv}}\) on a common measurable space \(\mathcal Z_{\mathrm{tv}}\), use
\[
\operatorname{TV}(\mu_{\mathrm{tv}},\nu_{\mathrm{tv}})
=
\sup_{E_{\mathrm{tv}}\text{ measurable in }\mathcal Z_{\mathrm{tv}}}
|\mu_{\mathrm{tv}}(E_{\mathrm{tv}})
 -\nu_{\mathrm{tv}}(E_{\mathrm{tv}})|.
\]
The conditional distribution of \((A,Y)\) under the cosine law is
\[
\begin{split}
\kappa_{\mathrm{cos},t_{\mathrm{cos}}}(x,\cdot)
={}&\frac{1+\tau_{P_{\mathrm{cos},t_{\mathrm{cos}}}}(x)}4
       \delta_{(1,1)}
+\frac{1-\tau_{P_{\mathrm{cos},t_{\mathrm{cos}}}}(x)}4
       \delta_{(1,-1)}\\
&+\frac14\delta_{(0,1)}+\frac14\delta_{(0,-1)}.
\end{split}
\]
For a measurable record event \(E_{\mathrm{tv}}\subseteq\mathcal O\), define its section by
\[
E_{\mathrm{tv},x}
=\{(a_{\mathrm{arm}},y):
  (x,a_{\mathrm{arm}},y)\in E_{\mathrm{tv}}\}.
\]
Subtracting the four displayed masses gives
\[
\begin{split}
&\kappa_{\mathrm{cos},0}(x,E_{\mathrm{tv},x})
-\kappa_{\mathrm{cos},t_{\mathrm{cos}}}(x,E_{\mathrm{tv},x})\\
&\quad=
-\frac{\tau_{P_{\mathrm{cos},t_{\mathrm{cos}}}}(x)}4
\left[
\mathbf1_{E_{\mathrm{tv},x}}(1,1)
-\mathbf1_{E_{\mathrm{tv},x}}(1,-1)
\right].
\end{split}
\]
The absolute difference of these two indicators is at most one. Therefore
\[
\left|
\kappa_{\mathrm{cos},0}(x,E_{\mathrm{tv},x})
-\kappa_{\mathrm{cos},t_{\mathrm{cos}}}(x,E_{\mathrm{tv},x})
\right|
\le
\frac{|\tau_{P_{\mathrm{cos},t_{\mathrm{cos}}}}(x)|}{4}
\le\frac{t_{\mathrm{cos}}}{16}.
\]
Integration over the common uniform design gives
\[
\begin{split}
|P_{\mathrm{cos},0}(E_{\mathrm{tv}})
-P_{\mathrm{cos},t_{\mathrm{cos}}}(E_{\mathrm{tv}})|
&\le
\int
\left|
\kappa_{\mathrm{cos},0}(x,E_{\mathrm{tv},x})
-\kappa_{\mathrm{cos},t_{\mathrm{cos}}}(x,E_{\mathrm{tv},x})
\right|\lambda(dx)\\
&\le\frac{t_{\mathrm{cos}}}{16},
\end{split}
\]
and hence
\[
\operatorname{TV}
(P_{\mathrm{cos},0},P_{\mathrm{cos},t_{\mathrm{cos}}})
\le\frac{t_{\mathrm{cos}}}{16}.
\]
% lean: CausalSmith.Stat.FinitepHomogeneityDensegamma.smallCosineLaw_tv

We use the expectation bound for any measurable
\(f_{\mathrm{tv}}:\mathcal Z_{\mathrm{tv}}\to[0,1]\). The layer-cake identity is
\[
\int f_{\mathrm{tv}}\,d\eta_{\mathrm{tv}}
=
\int_0^1
\eta_{\mathrm{tv}}
\{z_{\mathrm{tv}}:f_{\mathrm{tv}}(z_{\mathrm{tv}})\ge h_{\mathrm{tv}}\}
\,dh_{\mathrm{tv}},
\qquad
\eta_{\mathrm{tv}}\in\{\mu_{\mathrm{tv}},\nu_{\mathrm{tv}}\}.
\]
Subtracting these identities, taking absolute values, and bounding each event-probability difference gives
\[
\left|
\int f_{\mathrm{tv}}\,d\mu_{\mathrm{tv}}
-\int f_{\mathrm{tv}}\,d\nu_{\mathrm{tv}}
\right|
\le\operatorname{TV}(\mu_{\mathrm{tv}},\nu_{\mathrm{tv}}).
\]
Inserting the mixed product
\(P_{\mathrm{cos},0}\otimes
P_{\mathrm{cos},t_{\mathrm{cos}}}^{\otimes n_{\mathrm{prod}}}\)
between the two product laws, applying the triangle inequality, and integrating event sections gives
\[
\begin{split}
&\operatorname{TV}
\left(
P_{\mathrm{cos},0}^{\otimes(n_{\mathrm{prod}}+1)},
P_{\mathrm{cos},t_{\mathrm{cos}}}^{\otimes(n_{\mathrm{prod}}+1)}
\right)\\
&\quad\le
\operatorname{TV}(P_{\mathrm{cos},0},P_{\mathrm{cos},t_{\mathrm{cos}}})
+
\operatorname{TV}
\left(
P_{\mathrm{cos},0}^{\otimes n_{\mathrm{prod}}},
P_{\mathrm{cos},t_{\mathrm{cos}}}^{\otimes n_{\mathrm{prod}}}
\right),
\qquad n_{\mathrm{prod}}\in\mathbb N.
\end{split}
\]
For the first term, the integrated section probability is a \([0,1]\)-valued function, so the expectation bound applies; for the second, each section difference is bounded by the product total variation. At \(n_{\mathrm{prod}}=0\), both product laws are the point mass on the empty sample. Induction therefore yields
\[
\operatorname{TV}
\left(
P_{\mathrm{cos},0}^{\otimes n},
P_{\mathrm{cos},t_{\mathrm{cos}}}^{\otimes n}
\right)
\le
n\,\operatorname{TV}
(P_{\mathrm{cos},0},P_{\mathrm{cos},t_{\mathrm{cos}}})
\le\frac{nt_{\mathrm{cos}}}{16}.
\]
% lean: Causalean.Stat.Minimax.MomentMatchedMixture.tvDist_pi_iid_le

For the fixed \(n\ge2\), choose
\[
t_{\mathrm{small}}=\frac8{100n},
\qquad
0<t_{\mathrm{small}}\le1.
\]
The preceding product bound becomes
\[
\operatorname{TV}
\left(
P_{\mathrm{cos},0}^{\otimes n},
P_{\mathrm{cos},t_{\mathrm{small}}}^{\otimes n}
\right)
\le\frac{nt_{\mathrm{small}}}{16}
=\frac1{200}.
\]
For a binary-null level-valid test \(\phi\), define its seed average by
\[
\overline\phi(\boldsymbol o)
=\int_{[0,1]}\phi(\boldsymbol o,u_{\mathrm{seed}})
\,\lambda(du_{\mathrm{seed}}).
\]
It is measurable, satisfies \(0\le\overline\phi\le1\), and obeys
\[
\mathsf R_n(P,\phi)
=\int_{\mathcal O^n}\overline\phi\,dP^{\otimes n}.
\]
Applying the expectation bound to this average and using null validity gives
\[
\begin{split}
1-\mathsf R_n(P_{\mathrm{cos},t_{\mathrm{small}}},\phi)
&\ge
\frac9{10}
-\operatorname{TV}
\left(
P_{\mathrm{cos},0}^{\otimes n},
P_{\mathrm{cos},t_{\mathrm{small}}}^{\otimes n}
\right)\\
&\ge\frac{179}{200}>\frac1{10}.
\end{split}
\]
Whenever \(0<r\le t_{\mathrm{small}}d_0\), this cosine law belongs to the separated alternative set. Taking its error as a lower bound for the worst-case error, and then taking the infimum over level-valid tests, proves
\[
B_n^{\mathrm{bin}}(w,r)\ge\frac{179}{200}>\frac1{10}
\qquad(0<r\le t_{\mathrm{small}}d_0).
\]
% lean: CausalSmith.Stat.FinitepHomogeneityDensegamma.testingRiskOn_ge_of_two_point

Thus every successful separation in the binary capped-radius set exceeds \(t_{\mathrm{small}}d_0\). Its sentinel also satisfies
\[
t_{\mathrm{small}}d_0\le d_0\le D_w^{\mathrm{bin}}.
\]
The set is nonempty because it contains the sentinel. Taking its infimum, as defined in \cref{obj:def:binary-radius}, gives
\[
0<t_{\mathrm{small}}d_0
\le r_n^{*,\mathrm{bin}}(w).
\]
% lean: CausalSmith.Stat.FinitepHomogeneityDensegamma.binaryCriticalRadius_pos

\item Finally, the common cap and the risk equality give identical sets in the definitions of the binary and bounded radii:
\[
\begin{split}
&\left\{r\in(0,D_w^{\mathrm{bin}}):
 B_n^{\mathrm{bin}}(w,r)\le\frac1{10}\right\}
 \cup\{D_w^{\mathrm{bin}}\}\\
&\qquad=
\left\{r\in(0,D_w^{\mathrm b}):
 B_n^{\mathrm b}(w,r)\le\frac1{10}\right\}
 \cup\{D_w^{\mathrm b}\}.
\end{split}
\]
Their infima are therefore equal:
\[
r_n^{*,\mathrm{bin}}(w)=r_n^{*,\mathrm b}(w).
\]
The risk ordering also gives
\[
\begin{split}
&\left\{r\in(0,D_v):B_n(v,r)\le\frac1{10}\right\}\cup\{D_v\}\\
&\qquad\subseteq
\left\{r\in(0,D_v):B_n^{\mathrm b}(w,r)\le\frac1{10}\right\}
\cup\{D_v\}.
\end{split}
\]
Both sets contain \(D_v\) and are bounded below by zero, since \(D_v\ge d_0>0\). Taking infima reverses this inclusion. By
\cref{obj:def:bounded-radius,obj:def:critical-radius}, it follows that
\[
r_n^{*,\mathrm b}(w)\le r_n^*(v).
\]
Combining this with the positivity just proved establishes
\[
0<r_n^{*,\mathrm{bin}}(w)
=r_n^{*,\mathrm b}(w)\le r_n^*(v).
\]
The conversion properties and the identity for every \(n\in\mathbb N\) were established in the first step.
% lean: CausalSmith.Stat.FinitepHomogeneityDensegamma.cappedRadius_eq_of_risk_eq
% lean: CausalSmith.Stat.FinitepHomogeneityDensegamma.boundedCriticalRadius_le_criticalRadius
\end{enumerate}
\end{proof}

\begin{proof}[Proof of \cref{obj:thm:whole-null-calibration-resolved}]
Fix an admissible tuple \(v\) and an integer \(n\ge2\). Throughout, write
\[
M:=J_{\mathrm c},
\qquad
B:=\mathfrak b_{n,v}.
\]
These quantities have the values prescribed in \cref{obj:def:explicit-score-ledger}, including its small-sample branch.

\begin{enumerate}
\item First,
\[
B\ge0,
\qquad
M\ge1.
\]
For \(n<4\), these follow from \(B=0\) and \(M=1\). For \(n\ge4\), every summand in the formula for \(B\) is nonnegative: the clipping cutoffs and ranks are positive, and the increment weights \(a_j\) are nonnegative. This also covers an empty increment sum. The coarse rank is a nonnegative integer power of two, hence is at least one.
% lean: CausalSmith.Stat.FinitepHomogeneityDensegamma.ledgerBias_nonneg
% lean: CausalSmith.Stat.FinitepHomogeneityDensegamma.ledgerM_ge_one

\item Let \(P\in\mathcal M_v\). By \cref{obj:lem:original-score-covariance-ledger}, the affine coefficients in
\[
Z_b(c)=Z_{b,0}-cZ_{b,A},
\qquad b\in\{1,2\},
\]
have square-integrable scalar projections, and, for all \(u,z\in\mathbb R^M\),
\[
\begin{aligned}
\operatorname{Var}(u^T Z_{b,0})
&\le \Lambda_{n,v}\|u\|_2^2,\\
\operatorname{Var}(u^T Z_{b,A})
&\le \Lambda_{n,v}\|u\|_2^2,\\
\left|\operatorname{Cov}(u^T Z_{b,0},z^T Z_{b,A})\right|
&\le \Lambda_{n,v}\|u\|_2\|z\|_2.
\end{aligned}
\]
Taking the coordinate directions and summing their second moments proves square integrability of each coefficient vector.

Define the sample event
\[
\mathcal E_{\mathrm{prof}}(P)
:=
\left\{
\text{sample}:
\forall |c|\le\tfrac12,\ 
\left|W(c)-\|\theta_P(c)\|_2^2\right|
\le
128\left(
\sqrt{\Lambda_{n,v}}\|\theta_P(c)\|_2
+\sqrt M\,\Lambda_{n,v}
\right)
\right\}.
\]
By \cref{obj:lem:uniform-affine-profile-bound}, this event is measurable and
\[
P^{\otimes n}\!\left(\mathcal E_{\mathrm{prof}}(P)\right)\ge0.9925.
\]
The same result supplies, when \(\Lambda_{n,v}=0\), the simultaneous almost-sure identity
\[
W(c)=\|\theta_P(c)\|_2^2
\qquad\text{for every }|c|\le\tfrac12.
\]
% lean: CausalSmith.Stat.FinitepHomogeneityDensegamma.original_score_covariance_ledger
% lean: CausalSmith.Stat.FinitepHomogeneityDensegamma.uniform_affine_profile_bound

\item Suppose \(P\in H_0(v)\). By \cref{obj:def:null}, there is a constant \(c_0\) satisfying
\[
c_0\in[-\tfrac12,\tfrac12],
\qquad
\tau_P(x)=c_0\quad\text{for every }x\in[0,1].
\]
If \(n\ge4\), \cref{obj:lem:original-score-mean-ledger} gives
\[
\|\theta_P(c_0)\|_2\le B.
\]
If \(n<4\), the zero-score branch of \cref{obj:def:explicit-score-ledger} gives
\[
\theta_P(c)=0\quad\text{for every }c\in\mathbb R,
\qquad B=0,
\]
so the same bound holds. Since norms are nonnegative and \(c_0\) is an admissible candidate,
\[
\inf_{|c|\le1/2}\|\theta_P(c)\|_2
\le \|\theta_P(c_0)\|_2
\le B.
\]
% lean: CausalSmith.Stat.FinitepHomogeneityDensegamma.whole_null_calibration_resolved

\item We next establish the size bound for \(n\ge4\). The covariance ledger is strictly positive. Indeed, the construction gives
\[
s_n=\lfloor n/2\rfloor\ge2,
\qquad
T_0\ge1,
\qquad
V_0=32T_0^{2-p}>0,
\qquad
L_1\ge16\sqrt{V_0}>0.
\]
Consequently,
\[
\Lambda_{n,v}
=
\frac{L_1^2}{s_n}
+\frac{2L_2^2}{s_n(s_n-1)}
>0.
\]
% lean: CausalSmith.Stat.FinitepHomogeneityDensegamma.ledgerCov_pos

On \(\mathcal E_{\mathrm{prof}}(P)\), define
\[
u_{c_0}:=\|\theta_P(c_0)\|_2.
\]
Then \(0\le u_{c_0}\le B\), and the profiling inequality implies
\[
W(c_0)
\le
u_{c_0}^2
+128\sqrt{\Lambda_{n,v}}\,u_{c_0}
+128\sqrt M\,\Lambda_{n,v}.
\]
Completing the square gives
\[
0\le\left(B-64\sqrt{\Lambda_{n,v}}\right)^2,
\qquad
128\sqrt{\Lambda_{n,v}}\,B
\le B^2+4096\Lambda_{n,v}.
\]
Also, \(M\ge1\) implies
\[
\sqrt M\ge1,
\qquad
\Lambda_{n,v}\le\sqrt M\,\Lambda_{n,v}.
\]
Thus
\[
\begin{aligned}
W(c_0)
&\le B^2+128\sqrt{\Lambda_{n,v}}\,B
       +128\sqrt M\,\Lambda_{n,v}\\
&\le2B^2+4096\Lambda_{n,v}
       +128\sqrt M\,\Lambda_{n,v}\\
&\le2B^2+4224\sqrt M\,\Lambda_{n,v}\\
&\le2B^2+65536\sqrt M\,\Lambda_{n,v}.
\end{aligned}
\]
% lean: CausalSmith.Stat.FinitepHomogeneityDensegamma.profile_null_scalar

\item We justify the comparison with the computed profile minimum and the rounded cutoff. At a fixed sample, define
\[
Z_{b,0}:=Z_b(0),
\qquad
Z_{b,A}:=Z_b(0)-Z_b(1),
\qquad b\in\{1,2\}.
\]
The affine score construction gives \(Z_b(c)=Z_{b,0}-cZ_{b,A}\). Define its quadratic coefficients by
\[
\begin{aligned}
\omega_0&:=\langle Z_{1,0},Z_{2,0}\rangle,\\
\omega_1&:=-\langle Z_{1,A},Z_{2,0}\rangle
           -\langle Z_{1,0},Z_{2,A}\rangle,\\
\omega_2&:=\langle Z_{1,A},Z_{2,A}\rangle.
\end{aligned}
\]
Expansion yields
\[
W(c)=\omega_0+\omega_1c+\omega_2c^2.
\]
Define the smaller endpoint value and the computed profile value by
\[
\omega_{\mathrm{end}}
:=
\min\left\{
\omega_0-\frac{\omega_1}{2}+\frac{\omega_2}{4},
\omega_0+\frac{\omega_1}{2}+\frac{\omega_2}{4}
\right\},
\]
\[
\mathsf W_{\min}
:=
\begin{cases}
\displaystyle
\min\left\{\omega_{\mathrm{end}},
\omega_0-\frac{\omega_1^2}{4\omega_2}\right\},
&
\displaystyle
\omega_2>0,\quad
-\frac12\le-\frac{\omega_1}{2\omega_2}\le\frac12,\\[2mm]
\omega_{\mathrm{end}},&\text{otherwise}.
\end{cases}
\]
The quotient in the first branch is used under \(\omega_2>0\).

For every \(|c|\le1/2\),
\[
\mathsf W_{\min}\le W(c).
\]
To see this in the first branch, complete the square:
\[
W(c)-\left(\omega_0-\frac{\omega_1^2}{4\omega_2}\right)
=
\omega_2\left(c+\frac{\omega_1}{2\omega_2}\right)^2
\ge0.
\]
If \(\omega_2>0\) and the vertex lies to the left of the interval, then \(\omega_1>\omega_2\), and
\[
W(c)-W(-\tfrac12)
=
(c+\tfrac12)\bigl[\omega_1+\omega_2(c-\tfrac12)\bigr]
\ge0.
\]
If the vertex lies to the right, then \(\omega_1<-\omega_2\), and
\[
W(c)-W(\tfrac12)
=
(\tfrac12-c)\bigl[-\omega_1-\omega_2(c+\tfrac12)\bigr]
\ge0.
\]
Finally, if \(\omega_2\le0\), then \(c^2\le1/4\). For \(\omega_1\ge0\),
\[
W(c)-W(-\tfrac12)
=
\omega_1(c+\tfrac12)+\omega_2(c^2-\tfrac14)
\ge0,
\]
whereas for \(\omega_1<0\),
\[
W(c)-W(\tfrac12)
=
\omega_1(c-\tfrac12)+\omega_2(c^2-\tfrac14)
\ge0.
\]
These cases include a constant or linear statistic.
% lean: CausalSmith.Stat.FinitepHomogeneityDensegamma.ledger_profiledMin_le

For the cutoff, define
\[
x_{\mathrm{cut}}
:=2B^2+65536\sqrt M\,\Lambda_{n,v}>0,
\qquad
k_{\mathrm{cut}}
:=\left\lceil\log_2 x_{\mathrm{cut}}\right\rceil.
\]
The defining ceiling inequalities give
\[
\log_2 x_{\mathrm{cut}}
\le k_{\mathrm{cut}}
<\log_2 x_{\mathrm{cut}}+1.
\]
Exponentiating and using the cutoff formula therefore yields
\[
2B^2+65536\sqrt M\,\Lambda_{n,v}
\le h_{n,v}
\le4B^2+131072\sqrt M\,\Lambda_{n,v}.
\]
% lean: CausalSmith.Stat.FinitepHomogeneityDensegamma.ledgerCutoff_bounds

Combining these bounds with the preceding step, on \(\mathcal E_{\mathrm{prof}}(P)\),
\[
\mathsf W_{\min}
\le W(c_0)
\le2B^2+65536\sqrt M\,\Lambda_{n,v}
\le h_{n,v}.
\]
The strict rejection rule in \cref{obj:def:explicit-score-ledger} consequently gives
\[
\psi_{n,v}(\text{sample},\text{seed})=0
\quad\text{on }\mathcal E_{\mathrm{prof}}(P)\times[0,1].
\]

Define the experiment measure and its good event by
\[
\mathbb P_{\mathrm{exp}}:=P^{\otimes n}\otimes\lambda,
\qquad
\mathcal E_{\mathrm{exp}}(P)
:=\mathcal E_{\mathrm{prof}}(P)\times[0,1].
\]
Independence and the unit mass of \(\lambda\) imply
\[
\mathbb P_{\mathrm{exp}}\!\left(\mathcal E_{\mathrm{exp}}(P)\right)
=
P^{\otimes n}\!\left(\mathcal E_{\mathrm{prof}}(P)\right).
\]
The test takes values in \([0,1]\), so pointwise
\[
0\le\psi_{n,v}
\le\mathbf1_{\mathcal E_{\mathrm{exp}}(P)^c}.
\]
Integrating this inequality gives
\[
\begin{aligned}
\mathsf R_n(P,\psi_{n,v})
&\le
\mathbb P_{\mathrm{exp}}\!\left(\mathcal E_{\mathrm{exp}}(P)^c\right)\\
&=1-P^{\otimes n}\!\left(\mathcal E_{\mathrm{prof}}(P)\right)\\
&\le0.0075.
\end{aligned}
\]
For \(n<4\), the zero-test branch gives directly
\[
\mathsf R_n(P,\psi_{n,v})=0\le0.0075.
\]
% lean: CausalSmith.Stat.FinitepHomogeneityDensegamma.integral_le_bad_event
% lean: CausalSmith.Stat.FinitepHomogeneityDensegamma.whole_null_calibration_resolved

\item Now assume
\[
n\ge4,
\qquad
R_{\mathrm{att}}(n,v)<D_v,
\qquad
P\in\mathcal M_v,
\qquad
d(P)\ge R_{\mathrm{att}}(n,v).
\]
Work on \(\mathcal E_{\mathrm{prof}}(P)\). The computed profile value is attained: there exists
\[
c_{\min}\in[-\tfrac12,\tfrac12],
\qquad
W(c_{\min})=\mathsf W_{\min}.
\]
Indeed, the smaller endpoint value is attained at one of the two endpoints. In the vertex branch, choose that endpoint if its value is at most the vertex value; otherwise choose
\[
c_{\mathrm{vert}}:=-\frac{\omega_1}{2\omega_2},
\qquad
W(c_{\mathrm{vert}})
=\omega_0-\frac{\omega_1^2}{4\omega_2}.
\]
Here \(\omega_2>0\) and the branch condition places \(c_{\mathrm{vert}}\) in the candidate interval. In the remaining branch, choose the endpoint attaining \(\omega_{\mathrm{end}}\).
% lean: CausalSmith.Stat.FinitepHomogeneityDensegamma.ledger_profiledMin_attained

Define
\[
u_{\min}:=\|\theta_P(c_{\min})\|_2.
\]
By \cref{obj:lem:original-score-mean-ledger},
\[
u_{\min}
\ge\frac{3}{16}d(P)-16M^{-\gamma}-B.
\]
Substituting the assumed separation and its displayed definition gives
\[
\begin{aligned}
u_{\min}
&\ge
\frac{3}{16}R_{\mathrm{att}}(n,v)-16M^{-\gamma}-B\\
&=4B+1024M^{1/4}\sqrt{\Lambda_{n,v}}.
\end{aligned}
\]

We derive the scalar power bound. Since \(M\ge1\) and \(\Lambda_{n,v}>0\),
\[
M^{1/4}\ge1,
\qquad
\left(M^{1/4}\sqrt{\Lambda_{n,v}}\right)^2
=\sqrt M\,\Lambda_{n,v}>0.
\]
Thus
\[
u_{\min}\ge256\sqrt{\Lambda_{n,v}}\ge0,
\qquad
128\sqrt{\Lambda_{n,v}}\,u_{\min}
\le\frac12u_{\min}^2.
\]
Squaring the signal bound, with all its terms nonnegative, also gives
\[
\begin{aligned}
u_{\min}^2
&\ge
\left(4B+1024M^{1/4}\sqrt{\Lambda_{n,v}}\right)^2\\
&\ge16B^2+1048576\sqrt M\,\Lambda_{n,v}.
\end{aligned}
\]
The lower side of the profiling inequality now yields
\[
\begin{aligned}
W(c_{\min})
&\ge u_{\min}^2
      -128\sqrt{\Lambda_{n,v}}\,u_{\min}
      -128\sqrt M\,\Lambda_{n,v}\\
&\ge\frac12u_{\min}^2-128\sqrt M\,\Lambda_{n,v}\\
&\ge8B^2+524160\sqrt M\,\Lambda_{n,v}\\
&>4B^2+131072\sqrt M\,\Lambda_{n,v}.
\end{aligned}
\]
The last inequality is strict because \(\sqrt M\,\Lambda_{n,v}>0\).
% lean: CausalSmith.Stat.FinitepHomogeneityDensegamma.profile_alternative_scalar

Using the upper cutoff bound and attainment,
\[
h_{n,v}
\le4B^2+131072\sqrt M\,\Lambda_{n,v}
<W(c_{\min})
=\mathsf W_{\min}.
\]
Hence
\[
\psi_{n,v}=1
\quad\text{on }\mathcal E_{\mathrm{exp}}(P),
\qquad
0\le1-\psi_{n,v}
\le\mathbf1_{\mathcal E_{\mathrm{exp}}(P)^c}.
\]
Measurability and boundedness make the test integrable under the probability measure \(\mathbb P_{\mathrm{exp}}\). Integrating the last inequality gives
\[
\begin{aligned}
1-\mathsf R_n(P,\psi_{n,v})
&=\int(1-\psi_{n,v})\,d\mathbb P_{\mathrm{exp}}\\
&\le1-\mathbb P_{\mathrm{exp}}\!\left(\mathcal E_{\mathrm{exp}}(P)\right)\\
&=1-P^{\otimes n}\!\left(\mathcal E_{\mathrm{prof}}(P)\right)\\
&\le0.0075.
\end{aligned}
\]
% lean: CausalSmith.Stat.FinitepHomogeneityDensegamma.integral_le_bad_event
% lean: CausalSmith.Stat.FinitepHomogeneityDensegamma.whole_null_calibration_resolved

\item Finally, if \(n=2\) or \(n=3\), then \(n<4\), and the explicit small-sample branch gives
\[
\psi_{n,v}(\text{sample},\text{seed})=0
\qquad\text{for every sample and seed}.
\]
% lean: CausalSmith.Stat.FinitepHomogeneityDensegamma.ledger_small_sample
\end{enumerate}
\end{proof}

\begin{proof}[Proof of \cref{obj:thm:original-record-attainable-rate}]
\begin{enumerate}
\item Fix an admissible tuple \(v=(p,\alpha,\beta,\gamma)\). Its parameter restrictions give
\[
1<p\le2,\qquad
0<\alpha,\beta\le1,\qquad
\frac14\le\gamma\le1.
\]
Using the quantities in the statement, write
\[
D_p=1+2\alpha+\frac{\beta}{q}+\frac{S}{2\gamma}.
\]
Then
\[
0<q\le\frac12,\qquad S>0,\qquad t\ge0,\qquad D_p>1,
\qquad E_0>0,\qquad E_4>0.
\]
In particular,
\[
0<E\le\gamma\le1,\qquad
\frac{E}{S}<2,\qquad
\frac{E}{p-1}<1.
\]
Indeed, the three upper bounds follow from
\[
\frac{E}{\gamma}
\le\frac{2q}{2\gamma+q}\le1,
\qquad
\frac{E}{S}\le\frac{2}{D_p}<2,
\qquad
\frac{E}{p-1}
\le\frac{2\gamma}{p(2\gamma+q)}<1.
\]
The first uses \(q\le1/2\le2\gamma\), and the last uses \(p>1\) and \(q>0\).

For every positive argument of the geometric-series factor,
\[
0<2^{-x}<1,\qquad A(x)>0.
\]
Since \(\alpha>0\) and \(\lambda_0>0\), it follows that \(B_*\ge0\). Thus
\[
C^{\mathrm{att}}_v
=16\{16+5B_*+1024\sqrt{A_*}\}
\ge256>0.
\]
% lean: CausalSmith.Stat.FinitepHomogeneityDensegamma.CAtt_ge_256

\item The rule named \(\psi_{n,v}\) in
\cref{obj:thm:whole-null-calibration-resolved} is precisely
\(\phi^*_{n,v}\). That result supplies, for every integer \(n\ge2\),
\[
\mathsf R_n(P,\phi^*_{n,v})\le0.0075
\qquad(P\in H_0(v)).
\]
It also supplies power, for \(n\ge4\), at the calibrated separation
\[
R_{\mathrm{att}}(n,v)
=\frac{16}{3}
\left\{
16M^{-\gamma}+5B+1024M^{1/4}\sqrt{\Lambda_{n,v}}
\right\},
\qquad B=\mathfrak b_{n,v},
\]
provided \(R_{\mathrm{att}}(n,v)<D_v\).

The smallest sample sizes can be handled directly.
By \cref{obj:lem:causal-nonempty},
\[
d_0\le D_v\le\frac12.
\]
For \(n\in\{2,3\}\), the bound \(E\le1\) gives
\[
n^{-E}\ge n^{-1}\ge\frac13,
\qquad
D_v\le\frac12<\frac{256}{3}
\le C^{\mathrm{att}}_v n^{-E}.
\]
Consequently,
\[
C^{\mathrm{att}}_v n^{-E}<D_v
\quad\Longrightarrow\quad n\ge4
\qquad(n\ge2).
\]
The zero rule at \(n=2,3\) has the asserted size by the same calibration result.
We next establish the deterministic estimates needed to compare
\(R_{\mathrm{att}}(n,v)\) with the stated public level.
% lean: CausalSmith.Stat.FinitepHomogeneityDensegamma.attainment_small_sample_saturated

\item Let \(n\ge4\), and use the tuning of
\cref{obj:def:explicit-score-ledger,obj:def:blocks}. Explicitly,
\[
s=\lfloor n/2\rfloor\ge2,\qquad
a=s^{-E}\in(0,1],\qquad
M=\mathcal R(a^{-1/\gamma}),\qquad
K=\mathcal R\!\left(\max\{M,a^{-1/S}\}\right).
\]
Upper dyadic rounding obeys
\[
x_{\mathrm{dyad}}\le\mathcal R(x_{\mathrm{dyad}})
\le2x_{\mathrm{dyad}}
\qquad(x_{\mathrm{dyad}}\ge1),
\qquad
\mathcal R(2^j)=2^j\qquad(j\in\mathbb N_0),
\]
by the lower and upper ceiling bounds for \(\log_2 x_{\mathrm{dyad}}\).
It is also increasing. Hence
\[
a^{-1/\gamma}\le M\le2a^{-1/\gamma},
\qquad
K\ge a^{-1/S},
\]
and therefore
\[
M^{-\gamma}\le a,\qquad K^{-S}\le a.
\]

The bounds \(E/\gamma\le1\), \(E/S<2\), and \(E/(p-1)<1\) give
\[
a^{-1/\gamma}=s^{E/\gamma}\le s,
\qquad
a^{-1/S}=s^{E/S}\le s^2,
\qquad
1\le a^{-1/(p-1)}=s^{E/(p-1)}\le s.
\]
Thus \(M\le2s\). For the fine rank, if \(a^{-1/S}\le M\), dyadic idempotence gives \(K=M\le2s\le2s^2\); if \(M<a^{-1/S}\), rounding gives \(K\le2a^{-1/S}\le2s^2\). In either case,
\[
M\le2s,\qquad K\le2s^2.
\]

For the increment-cutoff bounds, the public quantities in either score branch are
\[
L=\log_2(K/M)\in\mathbb N_0,\qquad
R_j=2^jM\quad(0\le j\le L),
\]
\[
T_j=
\mathcal R\!\left(
\max\left\{
1,\min\left\{
a^{-1/(p-1)},
\left[
\frac{R_j^{-\alpha}}{a(R_j/K)^{\lambda_0}}
\right]^{1/(p-1)}
\right\}
\right\}
\right)
\quad(1\le j\le L).
\]
The single-cutoff score uses the main cutoff
\[
T_0=\mathcal R(a^{-1/(p-1)}).
\]
Since the argument defining each \(T_j\) lies between \(1\) and
\(a^{-1/(p-1)}\), rounding and monotonicity yield
\[
1\le T_0\le2s,\qquad
1\le T_j\le T_0\quad(1\le j\le L).
\]
These statements also cover \(L=0\), when the increment index set is empty.
% lean: CausalSmith.Stat.FinitepHomogeneityDensegamma.ledger_rate_bounds

\item We now bound the bias ledger. The main-cutoff estimate is
\[
T_0^{1-p}
\le\left(a^{-1/(p-1)}\right)^{1-p}
=a.
\]
In the single-cutoff branch, this and \(K^{-S}\le a\) immediately give
\[
B=400K^{-S}+20T_0^{1-p}\le420a\le B_*a.
\]

For the multilevel branch, \(S<\gamma\) and \(\alpha<q\). Define the positive local tail targets by
\[
z_{\mathrm{tail},j}
=\frac{R_j^{-\alpha}}{a(R_j/K)^{\lambda_0}}>0
\qquad(1\le j\le L).
\]
The cutoff formula gives
\[
T_j\ge
\min\left\{a^{-1/(p-1)},z_{\mathrm{tail},j}^{1/(p-1)}\right\}.
\]
If the minimum is the first term, its \((1-p)\)-power is \(a\); if it is the second term, that power is
\[
\left(z_{\mathrm{tail},j}^{1/(p-1)}\right)^{1-p}
=z_{\mathrm{tail},j}^{-1}
=a(R_j/K)^{\lambda_0}R_j^\alpha.
\]
Since \(1-p<0\), both cases imply
\[
T_j^{1-p}
\le a+a(R_j/K)^{\lambda_0}R_j^\alpha.
\]

The increment ranks satisfy \(M\ge1\) and \(R_L\le K\). Forward and backward geometric summation therefore give, for every \(r_{\mathrm{geom}}>0\),
\[
\begin{aligned}
\sum_{j=1}^{L}R_{j-1}^{-r_{\mathrm{geom}}}
&\le\sum_{j=0}^{L-1}2^{-jr_{\mathrm{geom}}}
\le A(r_{\mathrm{geom}}),\\
\sum_{j=1}^{L}(R_j/K)^{r_{\mathrm{geom}}}
&\le\sum_{j=0}^{L-1}2^{-jr_{\mathrm{geom}}}
\le A(r_{\mathrm{geom}}),\\
\sum_{j=1}^{L}R_j^{r_{\mathrm{geom}}}
&\le K^{r_{\mathrm{geom}}}A(r_{\mathrm{geom}}).
\end{aligned}
\]
For the second inequality, use \(R_j/K\le2^{-(L-j)}\); the third follows by factoring out \(K^{r_{\mathrm{geom}}}\). Each bound remains valid for an empty sum.

Using \(a_j=40R_{j-1}^{-\alpha}\), \(R_j=2R_{j-1}\), and \(2^\alpha\le2\), we obtain
\[
\begin{aligned}
a_jT_j^{1-p}
&\le40aR_{j-1}^{-\alpha}
 +40a\,2^\alpha(R_j/K)^{\lambda_0}\\
&\le80a\left\{R_{j-1}^{-\alpha}+(R_j/K)^{\lambda_0}\right\}.
\end{aligned}
\]
Summation consequently gives
\[
\sum_{j=1}^{L}a_jT_j^{1-p}
\le80a\{A(\alpha)+A(\lambda_0)\}.
\]
Substituting into the multilevel bias ledger proves
\[
\begin{aligned}
B
&=400K^{-S}+20T_0^{1-p}
  +10\sum_{j=1}^{L}a_jT_j^{1-p}\\
&\le\left[420+800\{A(\alpha)+A(\lambda_0)\}\right]a
=B_*a.
\end{aligned}
\]
Thus this bias estimate holds in both score branches.
% lean: CausalSmith.Stat.FinitepHomogeneityDensegamma.ledger_bias_bound

\item Consider first the single-cutoff covariance branch,
\[
S\ge\gamma\quad\text{or}\quad\alpha\ge q.
\]
If \(S\ge\gamma\), then \(E\le\gamma\le S\). If \(\alpha\ge q\), then
\[
E\le E_0\le q\le\alpha<S.
\]
Thus \(a^{-1/S}=s^{E/S}\le s\). The same two cases for fine-rank rounding used above now give
\[
K\le2s.
\]

Since \(0\le2-p\le1\) and \(T_0\le2a^{-1/(p-1)}\), the main second-moment quantity satisfies
\[
V_0=32T_0^{2-p}
\le32\,2^{2-p}a^{-t}
\le64a^{-t}.
\]
The single-cutoff formulas \(L_1=16\sqrt{V_0}\) and \(L_2=4\sqrt{KV_0}\) yield the exact identity
\[
\Lambda_{n,v}
=\frac{256V_0}{s}+\frac{32KV_0}{s(s-1)}.
\]
Because \(s\ge2\) and \(K\le2s\),
\[
32K\le64s\le128(s-1).
\]
It follows that
\[
\Lambda_{n,v}
\le\frac{384V_0}{s}
\le\frac{24576a^{-t}}{s}.
\]

The tail-channel identity and the choice \(E\le E_0\) give
\[
E_0\left(2+t+\frac{1}{2\gamma}\right)=1,
\qquad
E\left(2+t+\frac{1}{2\gamma}\right)\le1.
\]
Hence
\[
\frac{a^{-1/(2\gamma)}a^{-t}}{s}
=s^{E(1/(2\gamma)+t)-1}
\le s^{-2E}=a^2.
\]
Finally, \(M\le2a^{-1/\gamma}\) implies
\[
\begin{aligned}
\sqrt M\,\Lambda_{n,v}
&\le24576\sqrt2\,
 \frac{a^{-1/(2\gamma)}a^{-t}}{s}\\
&\le24576\sqrt2\,a^2
\le A_*a^2,
\end{aligned}
\]
where the last inequality follows from the displayed definition of \(A_*\).
% lean: CausalSmith.Stat.FinitepHomogeneityDensegamma.ledger_single_covariance_rate

\item Continue with the multilevel branch \(S<\gamma\), \(\alpha<q\).
Here \(a^{-1/\gamma}\le a^{-1/S}\). If \(a^{-1/S}\le M\), then
\[
K=M\le2a^{-1/\gamma}\le2a^{-1/S};
\]
if \(M<a^{-1/S}\), rounding gives the same final bound. Thus
\[
K\le2a^{-1/S},
\qquad
aK^\alpha
\le2^\alpha a^{1-\alpha/S}
=2^\alpha a^{\beta/S}\le2.
\]

The refinement depth satisfies
\[
2^L\le R_L\le K\le2a^{-1/S},
\qquad
L\log2\le\log2-\frac{\log a}{S}.
\]
The exponential tangent inequality, applied at \(-\log a-1\), gives
\[
-\log a
\le\exp(-\log a-1)=a^{-1}\exp(-1),
\qquad
-a\log a\le\exp(-1).
\]
Consequently,
\[
aL
\le a+\frac{-a\log a}{S\log2}
\le1+\frac{1}{\exp(1)S\log2}
=D_*.
\]

The bound \(T_j^{1-p}\le a+a(R_j/K)^{\lambda_0}R_j^\alpha\) can be rewritten using
\[
a(R_j/K)^{\lambda_0}R_j^\alpha
=aK^\alpha(R_j/K)^{\alpha+\lambda_0}.
\]
Geometric summation, \(aK^\alpha\le2\), and the fact that \(A\) decreases on positive arguments now give
\[
\begin{aligned}
\sum_{j=1}^{L}T_j^{1-p}
&\le aL+aK^\alpha A(\alpha+\lambda_0)\\
&\le D_*+2A(\alpha).
\end{aligned}
\]
The same geometric estimates supply
\[
\sum_{j=1}^{L}a_j\le80A(\alpha),
\]
and, from \(d_j=110R_{j-1}^{-s_0}+20T_j^{1-p}\),
\[
\sum_{j=1}^{L}d_j
\le110A(s_0)+20D_*+40A(\alpha).
\]
Moreover, \(T_j\le T_0\) and \(2-p\ge0\) imply
\[
\sqrt{V_j}\le\sqrt{V_0},
\qquad
\sum_{j=1}^{L}a_j\sqrt{V_j}
\le80A(\alpha)\sqrt{V_0}.
\]
Since \(T_0\ge1\), we also have \(\sqrt{V_0}\ge1\).
Inserting these bounds into the singleton ledger gives
\[
\begin{aligned}
L_1
&=16\sqrt{V_0}
 +8\sum_{j=1}^{L}(d_j+a_j\sqrt{V_j})\\
&\le
\left\{16+880A(s_0)+160D_*+960A(\alpha)\right\}
\sqrt{V_0}
=L_*\sqrt{V_0}.
\end{aligned}
\]
Using \(V_0\le64a^{-t}\), we obtain
\[
L_1^2\le64L_*^2a^{-t}.
\]
Together with \(a^{-1/(2\gamma)}a^{-t}/s\le a^2\), this proves
\[
\sqrt M\,\frac{L_1^2}{s}
\le\sqrt2\,64L_*^2
\frac{a^{-1/(2\gamma)}a^{-t}}{s}
\le\sqrt2\,64L_*^2a^2.
\]
% lean: CausalSmith.Stat.FinitepHomogeneityDensegamma.ledger_multilevel_singleton_covariance_rate

\item Still in the multilevel branch, define the backward summability exponent by
\[
\eta=1-(\alpha+\lambda_0)t.
\]
Because \(\alpha<q\) and \(t\ge0\),
\[
\eta\ge1-(q+\lambda_0)t
=\frac{(p-1)(p+6)}{4p}
=\eta_0>0.
\]
Also,
\[
0\le\alpha t=1-\eta-\lambda_0t<1.
\]

The positive-power cutoff estimate follows from
\[
T_j\le2\max\{1,z_{\mathrm{tail},j}^{1/(p-1)}\}.
\]
Indeed, using \(2^{2-p}\le2\) gives
\[
V_j=32T_j^{2-p}
\le64\max\{1,z_{\mathrm{tail},j}^{t}\},
\]
and therefore
\[
\sqrt{V_j}
\le\sqrt{64}\left(1+z_{\mathrm{tail},j}^{t/2}\right).
\]
Multiplying by \(\sqrt{R_j}\) and expanding the local target yields
\[
\sqrt{R_jV_j}
\le\sqrt{64}
\left\{
\sqrt{R_j}
+a^{-t/2}K^{\lambda_0t/2}R_j^{\eta/2}
\right\}.
\]
Backward geometric summation implies
\[
\sum_{j=1}^{L}\sqrt{R_j}
\le\sqrt K\,A(1/2),
\qquad
\sum_{j=1}^{L}R_j^{\eta/2}
\le K^{\eta/2}A(\eta/2)
\le K^{\eta/2}A(\eta_0/2).
\]
The fine-rank amplitude estimate also gives
\[
\begin{aligned}
\sqrt K
&=(aK^\alpha)^{t/2}
  a^{-t/2}K^{(1-\alpha t)/2}\\
&\le2^{t/2}a^{-t/2}K^{(1-\alpha t)/2}.
\end{aligned}
\]
Since \(\lambda_0t+\eta=1-\alpha t\), summing the guarded increment bound proves
\[
\begin{aligned}
\sum_{j=1}^{L}\sqrt{R_jV_j}
&\le\sqrt{64}
 \left\{2^{t/2}A(1/2)+A(\eta_0/2)\right\}
 a^{-t/2}K^{(1-\alpha t)/2}\\
&=G_*a^{-t/2}K^{(1-\alpha t)/2}.
\end{aligned}
\]
Now
\[
0<\frac{1-\alpha t}{2}\le\frac12,
\qquad
t+\frac{1-\alpha t}{S}=\frac{1+\beta t}{S}.
\]
Thus \(K\le2a^{-1/S}\) converts the last bound into
\[
\sum_{j=1}^{L}\sqrt{R_jV_j}
\le\sqrt2\,G_*a^{-(1+\beta t)/(2S)}.
\]

Using the definition of \(L_2\) and \((x+y)^2\le2x^2+2y^2\) for its two nonnegative summands, we get
\[
\begin{aligned}
L_2^2
&=16\left(
\sqrt{MV_0}+\sum_{j=1}^{L}\sqrt{R_jV_j}
\right)^2\\
&\le32MV_0
 +32\left(\sum_{j=1}^{L}\sqrt{R_jV_j}\right)^2\\
&\le4096s\,a^{-t}
 +64G_*^2a^{-(1+\beta t)/S}.
\end{aligned}
\]
Here the last line uses \(M\le2s\) and \(V_0\le64a^{-t}\).
Since \(s(s-1)\ge s^2/2\), it follows that
\[
\frac{2L_2^2}{s(s-1)}
\le\frac{16384a^{-t}}{s}
 +\frac{256G_*^2a^{-(1+\beta t)/S}}{s^2}.
\]

For the second noise channel, direct algebra gives
\[
2+t=\frac1q,
\qquad
2S+1+\beta t+\frac{S}{2\gamma}=D_p,
\]
and hence
\[
E_4\left(2+\frac{1+\beta t}{S}+\frac{1}{2\gamma}\right)=2.
\]
Using \(E\le E_4\), we obtain
\[
\frac{a^{-1/(2\gamma)}a^{-(1+\beta t)/S}}{s^2}
=s^{E(1/(2\gamma)+(1+\beta t)/S)-2}
\le s^{-2E}=a^2.
\]
Combining this with the tail-channel balance gives
\[
\sqrt M\,\frac{2L_2^2}{s(s-1)}
\le\sqrt2\{16384+256G_*^2\}a^2.
\]
Adding the bound \(\sqrt M\,L_1^2/s\le\sqrt2\,64L_*^2a^2\) therefore proves
\[
\begin{aligned}
\sqrt M\,\Lambda_{n,v}
&=\sqrt M\left\{\frac{L_1^2}{s}
                 +\frac{2L_2^2}{s(s-1)}\right\}\\
&\le\sqrt2\{64L_*^2+16384+256G_*^2\}a^2\\
&\le A_*a^2.
\end{aligned}
\]
This includes \(L=0\), since all the increment sums then vanish. Together with the single-cutoff bound \(\sqrt M\,\Lambda_{n,v}\le24576\sqrt2\,a^2\le A_*a^2\), the covariance estimate holds in both branches.
% lean: CausalSmith.Stat.FinitepHomogeneityDensegamma.ledger_multilevel_covariance_rate_of_singleton

\item The preceding estimates now give
\[
M^{-\gamma}\le a,\qquad
B\le B_*a,\qquad
M^{1/4}\sqrt{\Lambda_{n,v}}
=\sqrt{\sqrt M\,\Lambda_{n,v}}
\le\sqrt{A_*}\,a.
\]
Consequently,
\[
R_{\mathrm{att}}(n,v)
\le\frac{16}{3}
\{16+5B_*+1024\sqrt{A_*}\}\,a.
\]
The relation \(n\le3s\), together with \(0<E\le1\), implies
\[
a=s^{-E}\le3^E n^{-E}\le3n^{-E}.
\]
Thus
\[
R_{\mathrm{att}}(n,v)\le C^{\mathrm{att}}_v n^{-E}.
\]

If \(C^{\mathrm{att}}_v n^{-E}<D_v\), the bound \(D_v\le C^{\mathrm{att}}_v n^{-E}\) for \(n\in\{2,3\}\) ensures \(n\ge4\). The last inequality then gives
\[
R_{\mathrm{att}}(n,v)<D_v,\qquad
d(P)\ge C^{\mathrm{att}}_v n^{-E}
\ \Longrightarrow\
d(P)\ge R_{\mathrm{att}}(n,v).
\]
The power assertion of
\cref{obj:thm:whole-null-calibration-resolved} therefore yields
\[
1-\mathsf R_n(P,\phi^*_{n,v})\le0.0075
\]
for every stipulated alternative.

To obtain the capped-radius bound, observe that
\(C^{\mathrm{att}}_v n^{-E}>0\) and that both error bounds are below \(1/10\).
When \(C^{\mathrm{att}}_v n^{-E}<D_v\), the explicit test is admissible in
\cref{obj:def:testing-risk}, so
\[
B_n\!\left(v,C^{\mathrm{att}}_v n^{-E}\right)
\le0.0075\le\frac1{10}.
\]
For an empty alternative set its worst-error supremum has value zero, giving the same inequality. The defining set in
\cref{obj:def:critical-radius} therefore contains
\(C^{\mathrm{att}}_v n^{-E}\).
When \(D_v\le C^{\mathrm{att}}_v n^{-E}\), its sentinel \(D_v\) instead gives
\[
r_n^*(v)\le D_v\le C^{\mathrm{att}}_v n^{-E}.
\]
The defining set is nonempty and bounded below by zero in both cases. Hence
\[
r_n^*(v)\le C^{\mathrm{att}}_v n^{-E}
\qquad(n\ge2).
\]
% lean: CausalSmith.Stat.FinitepHomogeneityDensegamma.original_record_attainment_finite

\item Let \(n\ge N^{\mathrm{att}}_v\). By the ceiling definition,
\[
n\ge
\max\left\{4,
\left(\frac{2C^{\mathrm{att}}_v}{d_0}\right)^{1/E}\right\}.
\]
All bases are positive, and \(-E<0\). Taking this negative power therefore gives
\[
n^{-E}
\le
\left[
\left(\frac{2C^{\mathrm{att}}_v}{d_0}\right)^{1/E}
\right]^{-E}
=\left(\frac{2C^{\mathrm{att}}_v}{d_0}\right)^{-1}.
\]
Thus
\[
C^{\mathrm{att}}_v n^{-E}
\le\frac{d_0}{2}<d_0,
\]
which proves the witness-threshold assertion. The bounds \(M\le2s\), \(K\le2s^2\), \(1\le T_0\le2s\), and \(1\le T_j\le T_0\) for \(1\le j\le L\) give the tuning assertions.
% lean: CausalSmith.Stat.FinitepHomogeneityDensegamma.attainment_threshold

\item For compact uniformity, let
\[
\mathcal K\subseteq\mathcal V
\]
be compact in the product topology. At every admissible tuple,
\[
p-1>0,\quad S>0,\quad\gamma>0,\quad q>0,\quad D_p>0,
\quad s_0>0,\quad\lambda_0>0,\quad\eta_0>0.
\]
In particular, every geometric denominator appearing in the constants satisfies
\[
1-2^{-x}>0
\qquad
\left(
x\in
\{\alpha,\lambda_0,s_0,1/2,\eta_0/2\}
\right).
\]
The displayed formulas consequently show that \(E(v)\) and
\(C^{\mathrm{att}}_v\) are continuous at every admissible tuple: their operations are sums, products, quotients with nonzero denominators, minima, square roots, and real powers with positive bases.

Define the real threshold target on the admissible domain by
\[
\Xi(v)
=\max\left\{4,
\left(\frac{2C^{\mathrm{att}}_v}{d_0}\right)^{1/E(v)}
\right\}
\qquad(v\in\mathcal V).
\]
Since \(E(v)>0\), \(C^{\mathrm{att}}_v>0\), and \(d_0>0\), this function is continuous there as well.

If \(\mathcal K\) is nonempty, compactness supplies tuples
\(v_C,v_E,v_N\in\mathcal K\) satisfying
\[
C^{\mathrm{att}}_{v_C}
=\max_{v\in\mathcal K}C^{\mathrm{att}}_v,
\qquad
E(v_E)=\min_{v\in\mathcal K}E(v),
\qquad
\Xi(v_N)=\max_{v\in\mathcal K}\Xi(v).
\]
Choose
\[
C_{\mathcal K}=C^{\mathrm{att}}_{v_C},
\qquad
N_{\mathcal K}=\Xi(v_N)+1,
\qquad
e_{\mathcal K}=E(v_E)>0.
\]
The ceiling bound gives, for every \(v\in\mathcal K\),
\[
C^{\mathrm{att}}_v\le C_{\mathcal K},
\qquad
N^{\mathrm{att}}_v=\lceil\Xi(v)\rceil
<\Xi(v)+1\le N_{\mathcal K},
\qquad
e_{\mathcal K}\le E(v).
\]
If \(\mathcal K\) is empty, choose
\[
(C_{\mathcal K},N_{\mathcal K},e_{\mathcal K})=(0,0,1);
\]
the universal bounds are then vacuous.
% lean: CausalSmith.Stat.FinitepHomogeneityDensegamma.attainment_compact_uniformity

\item Finally, let \(w=(\alpha,\beta,\gamma)\in\mathcal W\), and define its matched admissible tuple by
\[
v_{\mathrm b}=(2,\alpha,\beta,\gamma)\in\mathcal V.
\]
At this tuple,
\[
q=\frac12,\qquad t=0,\qquad S=\alpha+\beta,
\]
so substitution into the two exponent formulas gives
\[
E_0(v_{\mathrm b})=\frac{2\gamma}{4\gamma+1},
\qquad
E_4(v_{\mathrm b})
=\frac{2S}{1+2S+S/(2\gamma)}.
\]
Taking their minimum proves
\[
E(v_{\mathrm b})=E_{\mathrm b}.
\]

The definitions in
\cref{obj:def:bounded-model,obj:def:bounded-null,obj:def:sharp-frontier-scales}
give
\[
\mathcal M_w^{\mathrm b}\subseteq\mathcal M_{v_{\mathrm b}},
\qquad
H_0^{\mathrm b}(w)\subseteq H_0(v_{\mathrm b}),
\qquad
\phi^{*,\mathrm b}_{n,w}=\phi^*_{n,v_{\mathrm b}},
\]
and
\[
C^{\mathrm{att}}_{(2,w)}n^{-E_{\mathrm b}}
=C^{\mathrm{att}}_{v_{\mathrm b}}n^{-E(v_{\mathrm b})}>0.
\]
Moreover, \cref{obj:prop:bounded-submodel-comparison} supplies the exact cap identity
\[
D_w^{\mathrm b}=D_{v_{\mathrm b}}\ge d_0.
\]
The original-model size bound therefore applies to every bounded null law. If
\[
C^{\mathrm{att}}_{(2,w)}n^{-E_{\mathrm b}}<D_w^{\mathrm b},
\]
the cap identity makes the original-model power condition hold at
\(v_{\mathrm b}\). Thus every bounded-model law with distance at least this level satisfies
\[
1-\mathsf R_n(P,\phi^{*,\mathrm b}_{n,w})\le0.0075.
\]

When this positive level is below \(D_w^{\mathrm b}\), the resulting bounded-model risk bound is
\[
B_n^{\mathrm b}
\!\left(w,C^{\mathrm{att}}_{(2,w)}n^{-E_{\mathrm b}}\right)
\le0.0075\le\frac1{10},
\]
so the level belongs to the defining set of
\cref{obj:def:bounded-radius}. When it is at least \(D_w^{\mathrm b}\), the sentinel gives the radius bound directly. Therefore, for every \(n\ge2\),
\[
r_n^{*,\mathrm b}(w)
\le C^{\mathrm{att}}_{(2,w)}n^{-E_{\mathrm b}}.
\]
% lean: CausalSmith.Stat.FinitepHomogeneityDensegamma.bounded_record_attainment_finite
\end{enumerate}
\end{proof}

\begin{proof}[Proof of \cref{obj:thm:tail-essential-full-record-converse}]
Fix an admissible tuple \(v=(p,\alpha,\beta,\gamma)\), and put
\(w=(\alpha,\beta,\gamma)\). Throughout, use the quantities and selected
constructions of
\cref{obj:def:sharp-frontier-scales,obj:def:converse-handle}.

\begin{enumerate}
\item
Fix \(n\ge2\), and first suppose \(F(p)<1\).
By \cref{obj:lem:rough-full-record-sharp-lower}, both selected priors are
normalized, and
\[
\begin{gathered}
\operatorname{supp}_+(\pi_{0,n,v})\subseteq H_0(v),\\
\operatorname{supp}_+(\pi_{1,n,v})
\subseteq
\{P\in\mathcal M_v:c_vn^{-E_4(p)}\le d(P)<d_0\},\\
d_0\le D_v,\qquad
\operatorname{TV}(\mathbb P_{0,n,v},\mathbb P_{1,n,v})<\frac12,\qquad
B_n(v,c_vn^{-E_4(p)})\ge\frac25.
\end{gathered}
\]
These conclusions cover every \(n\ge2\), including the paired-tent
selection below \(N_{\mathrm{low}}(v)\).

The scale identification follows by subtracting the two exponents:
\[
\begin{aligned}
E_4(p)-E_0(p)
&=
\frac{2S(2\gamma+q)-2\gamma qD_p}
     {(2\gamma+q)D_p}\\
&=
\frac{2\gamma q}{(2\gamma+q)D_p}\{F(p)-1\}.
\end{aligned}
\]
Admissibility gives \(q>0\), \(S>0\), \(\gamma>0\), and \(D_p>0\).
Thus \(F(p)<1\) implies
\[
E(p)=E_4(p),\qquad \rho_n(v)=n^{-E_4(p)}.
\]

We also verify the arm support at the selected magnitude. When
\(n\ge N_{\mathrm{low}}(v)\), the rank conditions are supplied by
\cref{obj:lem:rough-full-record-sharp-lower}, and the tuning agrees with
\cref{obj:lem:copula-frame-legality}:
\[
b=\frac{K^{-\beta}}{256},\qquad
\varepsilon_{n,v}=(16b)^{1/q}
=16^{-1/q}K^{-\beta/q},\qquad
L_{n,v}=\varepsilon_{n,v}^{-1/p}.
\]
The normalized table in \cref{obj:def:marked-table} assigns its outcome
mass to \(0,-L_{n,v},L_{n,v}\). Conditioning on either treatment arm
therefore retains all its mass on those three points. When
\(n<N_{\mathrm{low}}(v)\), the selected paired-tent construction has the
same arm-support conclusion by
\cref{obj:lem:paired-tent-full-record-lower}. Consequently, in either
selection,
\[
Q_{a,P}\bigl(\{0,-L_{n,v},L_{n,v}\}\mid x\bigr)=1
\]
for each \(a\in\{0,1\}\), for \(\lambda\)-almost every \(x\), and for every
positive-weight law in either prior.

Finally, normalization and nonnegative weights ensure that the
alternative prior has a positive-weight law. Choosing such a law \(P\)
gives
\[
c_v\rho_n(v)\le d(P)<d_0\le D_v.
\]
Moreover, \(N_{\mathrm{low}}(v)\ge2\), \(k_v>0\), and
\(c_0^{\mathrm e}>0\), so the defining minimum for \(c_v\) is positive.
Since \(n>0\), also \(\rho_n(v)>0\). Hence
\[
0<c_v\rho_n(v)<D_v,
\]
and every positive-weight alternative law has distance strictly below
\(D_v\). This proves the first group of conclusions in the rough phase.
% lean: CausalSmith.Stat.FinitepHomogeneityDensegamma.rough_phase_converse_assembly

\item
Suppose instead \(F(p)\ge1\). The exponent-difference identity from the
preceding step gives
\[
E(p)=E_0(p),\qquad
c_v=c_0^{\mathrm e},\qquad
c_v\rho_n(v)=c_0^{\mathrm e}n^{-E_0(p)}.
\]
The selector of \cref{obj:def:converse-handle} chooses the paired-tent
priors for every \(n\ge2\). Their finite sign distributions are
normalized. All null sign configurations give the same zero-effect
null law \(P_0\), so
\[
\pi_{0,n,v}=\delta_{P_0},\qquad
\mathbb P_{0,n,v}=P_0^{\otimes n}.
\]
By \cref{obj:lem:paired-tent-full-record-lower},
\[
\begin{gathered}
P_0\in H_0(v),\\
\operatorname{supp}_+(\pi_{1,n,v})
\subseteq
\{P\in\mathcal M_v:
c_0^{\mathrm e}n^{-E_0(p)}\le d(P)<d_0\},\\
d_0\le D_v,\qquad
\operatorname{TV}(\mathbb P_{0,n,v},\mathbb P_{1,n,v})<\frac12,\\
B_n(v,c_0^{\mathrm e}n^{-E_0(p)})\ge\frac25.
\end{gathered}
\]
That construction also gives the required three-point conditional arm
support at \(L_{n,v}\). Its alternative prior has a positive-weight
law; for any such law,
\[
c_v\rho_n(v)\le d(P)<d_0\le D_v.
\]
Positivity of \(c_0^{\mathrm e}\) and \(n^{-E_0(p)}\) yields
\(0<c_v\rho_n(v)<D_v\). Substituting the displayed scale identity proves
the first group of conclusions in this phase as well. The equality
case \(F(p)=1\) is included here.
% lean: CausalSmith.Stat.FinitepHomogeneityDensegamma.nonrough_phase_converse_assembly

\item
We next prove divergence of the selected magnitude, again splitting at
\(F(p)<1\). In that phase, for all
\(n\ge N_{\mathrm{low}}(v)\), the selected fine rank satisfies
\[
K\ge H_{\mathrm{fine}}n^{2/D_p}\longrightarrow\infty,
\]
and moment normalization gives
\[
L_{n,v}
=
\left(16^{-1/q}K^{-\beta/q}\right)^{-1/p}
=
16^{1/(pq)}K^{\beta/(pq)}.
\]
Both \(2/D_p\) and \(\beta/(pq)\) are positive, so
\(L_{n,v}\to\infty\). The finitely many preceding selections leave this
limit unchanged.

When \(F(p)\ge1\), the paired-tent formula applies throughout:
\[
N=2\left\lceil n^{2q/(2\gamma+q)}\right\rceil
\ge n^{2q/(2\gamma+q)}\longrightarrow\infty,
\qquad
L_{n,v}=N^{\gamma/(p-1)}.
\]
Here \(2q/(2\gamma+q)>0\) and \(\gamma/(p-1)>0\), proving the same
divergence.
% lean: CausalSmith.Stat.FinitepHomogeneityDensegamma.magnitude_tendsto_atTop

\item
Assume \(p<2\). At the matched smoothness tuple, write
\[
D_2=1+2S+\frac{S}{2\gamma}.
\]
The identity \(q^{-1}=2+(2-p)/(p-1)\) gives
\[
D_p=D_2+\frac{\beta(2-p)}{p-1}.
\]
Direct subtraction then yields the two strictly positive differences
\[
E_0(2)-E_0(p)
=
\frac{4\gamma^2(2-p)}
     {(4\gamma+1)(2\gamma p+p-1)}>0,
\]
\[
E_4(2)-E_4(p)
=
\frac{2S\beta(2-p)}
     {(p-1)D_pD_2}>0.
\]
All denominators are positive by admissibility. Since \(E(p)\) is the
minimum of its two component exponents,
\[
E(p)\le E_0(p)<E_0(2),\qquad
E(p)\le E_4(p)<E_4(2).
\]
Thus the exponent gap, defined on this face by
\[
\delta_v=E(2)-E(p),
\]
is strictly positive. For every positive integer \(n\),
\[
\frac{\rho_n(v)}{\rho_n^{\mathrm b}(w)}
=n^{\delta_v}\longrightarrow\infty.
\]
% lean: CausalSmith.Stat.FinitepHomogeneityDensegamma.matched_scale_ratio_tendsto

\item
Continue with \(p<2\). Since \(c_v>0\), the preceding limit permits an
integer \(N_{\mathrm{dom}}(v)\in\mathbb N\) such that
\[
\frac{\rho_n(v)}{\rho_n^{\mathrm b}(w)}
\ge\frac{C_w^{\mathrm b}}{c_v}
\qquad
\text{for every integer }n\ge N_{\mathrm{dom}}(v)\text{ with }n>0.
\]
Multiplying by the positive bounded scale and by \(c_v\) gives
\[
C_w^{\mathrm b}\rho_n^{\mathrm b}(w)\le c_v\rho_n(v).
\]
Define the required threshold by
\[
N(v)=\max\{2,N_{\mathrm{dom}}(v),N_w^{\mathrm b}\}.
\]
For \(n\ge N(v)\), \cref{obj:thm:original-record-attainable-rate}, applied
at \((2,w)\), and \cref{obj:prop:bounded-submodel-comparison} give
\[
C_w^{\mathrm b}\rho_n^{\mathrm b}(w)<d_0\le D_w^{\mathrm b}.
\]
The bounded total test therefore satisfies
\[
\mathsf R_n(P,\phi^{*,\mathrm b}_{n,w})\le0.0075
\quad(P\in H_0^{\mathrm b}(w)),
\]
and
\[
1-\mathsf R_n(P,\phi^{*,\mathrm b}_{n,w})\le0.0075
\quad
(P\in\mathcal M_w^{\mathrm b},\ d(P)\ge c_v\rho_n(v)).
\]
Its size is below \(1/10\), so
\[
B_n^{\mathrm b}(w,c_v\rho_n(v))\le0.0075.
\]

Let \(\pi_0,\pi_1\) be any normalized finite priors satisfying the bounded
support conditions in the statement. We establish the two-prior
inequality used to compare their mixtures. For any randomized test
\(\phi\), define its seed average on the complete record space by
\[
f_\phi:
\bigl([0,1]\times\{0,1\}\times\mathbb R\bigr)^n\longrightarrow[0,1],
\qquad
f_\phi(\mathcal D_n)=\int_{[0,1]}\phi(\mathcal D_n,U)\,\lambda(dU).
\]
Bounded product integration and finite-sum linearity give, for either
prior,
\[
\int f_\phi\,d\mathbb P_{\pi_\nu,n}
=
\int \mathsf R_n(P,\phi)\,\pi_\nu(dP),
\qquad \nu\in\{0,1\}.
\]
The bounded-function form of the total-variation inequality gives
\[
\left|
\int f_\phi\,d\mathbb P_{\pi_1,n}
-
\int f_\phi\,d\mathbb P_{\pi_0,n}
\right|
\le
\operatorname{TV}(\mathbb P_{\pi_0,n},\mathbb P_{\pi_1,n}).
\]
If \(\phi\) has size at most \(1/10\) on the bounded null, normalization
and the null support condition imply
\[
\int \mathsf R_n(P,\phi)\,\pi_0(dP)\le\frac1{10}.
\]
The alternative support condition and normalization consequently yield
\[
\begin{aligned}
\sup_{\substack{P\in\mathcal M_w^{\mathrm b}\\
                 d(P)\ge c_v\rho_n(v)}}
\{1-\mathsf R_n(P,\phi)\}
&\ge
\int \{1-\mathsf R_n(P,\phi)\}\,\pi_1(dP)\\
&\ge
\frac9{10}
-\operatorname{TV}(\mathbb P_{\pi_0,n},\mathbb P_{\pi_1,n}).
\end{aligned}
\]
Weights equal to zero contribute zero in these averages; every
positive-weight law satisfies the stipulated support condition.
Normalization ensures a positive-weight alternative law, and the zero
rejection map is an admissible test. Taking the infimum over
null-sized tests therefore gives
\[
\frac9{10}
-\operatorname{TV}(\mathbb P_{\pi_0,n},\mathbb P_{\pi_1,n})
\le B_n^{\mathrm b}(w,c_v\rho_n(v)).
\]
Combining this with the attained risk bound proves
\[
\operatorname{TV}(\mathbb P_{\pi_0,n},\mathbb P_{\pi_1,n})
\ge0.8925\ge\frac12.
\]
The threshold depends on \(v\), and the argument applies to every
normalized finite prior pair in the statement.
% lean: CausalSmith.Stat.FinitepHomogeneityDensegamma.bounded_tail_essentiality

\item
Now assume \(p=2\), and fix \(n\ge2\). Then
\(v=(2,\alpha,\beta,\gamma)\).
We first verify the original-model support conclusions for the selected
bounded priors, splitting at \(F(2)<1\).

In the rough phase, \cref{obj:lem:rough-full-record-sharp-lower} supplies
\[
\begin{gathered}
\operatorname{supp}_+(\pi^{\mathrm b}_{0,n,v})
\subseteq H_0^{\mathrm b}(w),\\
\operatorname{supp}_+(\pi^{\mathrm b}_{1,n,v})
\subseteq
\{P\in\mathcal M_w^{\mathrm b}:
c_vn^{-E_4(2)}\le d(P)<d_0\}.
\end{gathered}
\]
By the phase identity, \(\rho_n(v)=n^{-E_4(2)}\).
The model inclusion of \cref{obj:prop:bounded-submodel-comparison}
preserves the effect, and the null constancy condition is the same in
both models. Hence
\[
\begin{gathered}
\operatorname{supp}_+(\pi^{\mathrm b}_{0,n,v})\subseteq H_0(v),\\
\operatorname{supp}_+(\pi^{\mathrm b}_{1,n,v})
\subseteq
\{P\in\mathcal M_v:c_v\rho_n(v)\le d(P)\}.
\end{gathered}
\]

In the complementary phase \(F(2)\ge1\), the selected bounded priors are
the signed-binary paired-tent priors. To identify their null support,
define the balanced deterministic zero-effect law by
\[
P_{\mathrm{zero}}(dx,dA,dy)
=
\lambda(dx)\,
\frac{\delta_0+\delta_1}{2}(dA)\,
\delta_0(dy).
\]
It has propensity \(1/2\), baseline mean zero, and effect zero. Its
smoothness conditions hold because these functions are constant, and
its raw moment is zero. Its signed-binary conversion from
\cref{obj:prop:bounded-submodel-comparison} is
\[
P_{\mathrm{zero}}^{\mathrm{bin}}(dx,dA,dy)
=
\lambda(dx)\,
\frac{\delta_0+\delta_1}{2}(dA)\,
\frac{\delta_{-1}+\delta_1}{2}(dy).
\]
This conversion preserves the means and effect, and its conditional
second moment is one. Thus
\(P_{\mathrm{zero}}^{\mathrm{bin}}\in H_0(v)\).
Every selected null sign configuration gives precisely this law, so
\[
\operatorname{supp}_+(\pi^{\mathrm b}_{0,n,v})
\subseteq\{P_{\mathrm{zero}}^{\mathrm{bin}}\}\subseteq H_0(v).
\]
For the alternatives,
\cref{obj:lem:paired-tent-full-record-lower} and
\cref{obj:prop:bounded-submodel-comparison} give original-model
membership and
\[
d(P)\ge c_0^{\mathrm e}n^{-E_0(2)}
=c_v\rho_n(v)
\qquad
(P\in\operatorname{supp}_+(\pi^{\mathrm b}_{1,n,v})).
\]
This proves the original-model support conclusions in both phases.
% lean: CausalSmith.Stat.FinitepHomogeneityDensegamma.second_moment_rough_prior_support

\item
The bounded risk lower bound follows from the same phase split.
When \(F(2)<1\),
\cref{obj:lem:rough-full-record-sharp-lower} gives
\[
B_n^{\mathrm b}(w,c_vn^{-E_4(2)})\ge\frac25.
\]
When \(F(2)\ge1\),
\cref{obj:lem:paired-tent-full-record-lower} gives
\[
B_n^{\mathrm b}(w,c_0^{\mathrm e}n^{-E_0(2)})\ge\frac25.
\]
The phase identities identify each displayed separation with
\(c_v\rho_n(v)\). The risk comparison in
\cref{obj:prop:bounded-submodel-comparison} consequently yields
\[
\frac25
\le B_n^{\mathrm b}(w,c_v\rho_n(v))
\le B_n(v,c_v\rho_n(v)).
\]
% lean: CausalSmith.Stat.FinitepHomogeneityDensegamma.second_moment_testingRisk_lower

\item
We assemble the remaining second-moment conclusions. Evaluation at
\(v=(2,w)\), together with
\cref{obj:prop:bounded-submodel-comparison}, gives
\[
\begin{gathered}
\rho_n(v)=\rho_n^{\mathrm b}(w),\qquad
c_v=c_w^{\mathrm b},\qquad
C_v=C_w^{\mathrm b},\qquad
N_v=N_w^{\mathrm b},\\
D_v=D_w^{\mathrm b}.
\end{gathered}
\]
For \(F(2)<1\), the selected-prior conclusions of
\cref{obj:lem:rough-full-record-sharp-lower} apply for every \(n\ge2\),
covering the bounded rough construction above its threshold and the
signed-binary paired-tent construction below it. For \(F(2)\ge1\), use
the signed-binary conclusions of
\cref{obj:lem:paired-tent-full-record-lower}. The signed-binary model
inclusion of \cref{obj:prop:bounded-submodel-comparison} supplies bounded
membership, with null constancy retained. In either case, both selected
bounded priors are normalized and
\[
\begin{gathered}
\operatorname{supp}_+(\pi^{\mathrm b}_{0,n,v})
\subseteq H_0^{\mathrm b}(w),\\
\operatorname{supp}_+(\pi^{\mathrm b}_{1,n,v})
\subseteq
\{P\in\mathcal M_w^{\mathrm b}:
c_v\rho_n(v)\le d(P)<d_0\},\\
\operatorname{TV}(\mathbb P^{\mathrm b}_{0,n,v},
                  \mathbb P^{\mathrm b}_{1,n,v})<\frac12,\qquad
B_n^{\mathrm b}(w,c_v\rho_n(v))\ge\frac25.
\end{gathered}
\]
For any finite representation of either prior, nonnegative weights
summing to one satisfy
\[
\sum_{j:\,\omega_j>0}\omega_j=\sum_j\omega_j=1.
\]
Indeed, every omitted weight is zero. At least one weight is positive,
and combining weights of coincident laws preserves their sum.
Therefore
\[
\operatorname{supp}_+(\pi^{\mathrm b}_{\nu,n,v})\ne\varnothing,
\qquad
\sum_{P\in\operatorname{supp}_+(\pi^{\mathrm b}_{\nu,n,v})}
\pi^{\mathrm b}_{\nu,n,v}(\{P\})=1,
\qquad \nu\in\{0,1\}.
\]
Choose a positive-weight alternative law \(P\). Its placement, and the
witness lower bound from
\cref{obj:prop:bounded-submodel-comparison}, give
\[
0<c_v\rho_n(v)\le d(P)<d_0\le D_w^{\mathrm b}.
\]
% lean: CausalSmith.Stat.FinitepHomogeneityDensegamma.second_moment_bounded_receipt

\item
We obtain the lower capped-radius bound first in the bounded model.
Every successful bounded separation satisfies
\[
r_{\mathrm{succ}}\in(0,D_w^{\mathrm b}),
\qquad
B_n^{\mathrm b}(w,r_{\mathrm{succ}})\le\frac1{10}.
\]
If \(r_{\mathrm{succ}}<c_v\rho_n(v)\), containment of the separated
alternative classes gives
\[
B_n^{\mathrm b}(w,c_v\rho_n(v))
\le B_n^{\mathrm b}(w,r_{\mathrm{succ}})
\le\frac1{10},
\]
contradicting the lower bound \(2/5\). The cap
\(D_w^{\mathrm b}\) also exceeds \(c_v\rho_n(v)\). Thus every member of
the nonempty set defining the capped infimum is at least
\(c_v\rho_n(v)\), and
\[
c_v\rho_n(v)\le r_n^{*,\mathrm b}(w).
\]
Now \cref{obj:prop:bounded-submodel-comparison} gives
\[
c_v\rho_n(v)
\le r_n^{*,\mathrm b}(w)
=r_n^{*,\mathrm{bin}}(w)
\le r_n^*(v).
\]
The upper capped-radius bounds and both size bounds follow from
\cref{obj:thm:original-record-attainable-rate}:
\[
r_n^*(v)\le C_v\rho_n(v),\qquad
r_n^{*,\mathrm b}(w)\le C_w^{\mathrm b}\rho_n^{\mathrm b}(w),
\]
\[
\mathsf R_n(P,\phi^*_{n,v})\le0.0075
\quad(P\in H_0(v)),\qquad
\mathsf R_n(P,\phi^{*,\mathrm b}_{n,w})\le0.0075
\quad(P\in H_0^{\mathrm b}(w)).
\]
Substituting the matching scale and multiplier identities proves both
asserted two-sided radius bounds.
% lean: CausalSmith.Stat.FinitepHomogeneityDensegamma.second_moment_capped_lower

\item
Suppose \(C_v\rho_n(v)<D_v\). The matching identities give
\[
C_w^{\mathrm b}\rho_n^{\mathrm b}(w)<D_w^{\mathrm b}.
\]
To prove nonemptiness in the bounded model, suppose its separated class
were empty. Every bounded-model law would then satisfy
\[
d(P)<C_w^{\mathrm b}\rho_n^{\mathrm b}(w).
\]
The bounded model is nonempty, as witnessed by a positive-weight law
from the selected alternative prior. The displayed bound would imply
\[
D_w^{\mathrm b}
\le C_w^{\mathrm b}\rho_n^{\mathrm b}(w),
\]
contradicting the strict inequality. Hence there is a bounded-model law
with
\[
d(P)\ge C_w^{\mathrm b}\rho_n^{\mathrm b}(w)=C_v\rho_n(v).
\]
Its original-model membership follows from
\cref{obj:prop:bounded-submodel-comparison}, proving nonemptiness of both
separated classes.

Under this same strict-separation condition,
\cref{obj:thm:original-record-attainable-rate} supplies
\[
1-\mathsf R_n(P,\phi^*_{n,v})\le0.0075
\quad(P\in\mathcal M_v,\ d(P)\ge C_v\rho_n(v)),
\]
and
\[
1-\mathsf R_n(P,\phi^{*,\mathrm b}_{n,w})\le0.0075
\quad
(P\in\mathcal M_w^{\mathrm b},\
d(P)\ge C_w^{\mathrm b}\rho_n^{\mathrm b}(w)).
\]
% lean: CausalSmith.Stat.FinitepHomogeneityDensegamma.second_moment_face_package

\item
Finally, let \(n\ge N_v\). Positivity of the component exponents implies
\(E(2)>0\), and
\cref{obj:thm:original-record-attainable-rate} gives \(C_v>0\).
The threshold formula yields
\[
n\ge
\max\left\{4,\left(\frac{2C_v}{d_0}\right)^{1/E(2)}\right\}.
\]
Since raising positive numbers to the power \(-E(2)\) reverses their
order,
\[
\begin{aligned}
C_v\rho_n(v)
&=C_vn^{-E(2)}\\
&\le
C_v\left\{
\left(\frac{2C_v}{d_0}\right)^{1/E(2)}
\right\}^{-E(2)}
=\frac{d_0}{2}.
\end{aligned}
\]
Also, the witness-threshold conclusion of
\cref{obj:thm:original-record-attainable-rate} gives
\[
C_v\rho_n(v)<d_0\le D_v.
\]
The matching identities therefore imply
\[
C_w^{\mathrm b}\rho_n^{\mathrm b}(w)<D_w^{\mathrm b}.
\]
At \(n=2,3\), the small-sample branch of
\cref{obj:def:explicit-score-ledger} sets the total rule to zero.
Evaluating that same branch at \(v=(2,w)\) gives, for every sample \(z\),
\[
\phi^*_{n,v}(z)=0,\qquad
\phi^{*,\mathrm b}_{n,w}(z)=0.
\]
This completes all conclusions on the second-moment face.
% lean: CausalSmith.Stat.FinitepHomogeneityDensegamma.second_moment_attainment_half
\end{enumerate}
\end{proof}

\begin{proof}[Proof of \cref{obj:thm:heavy-tail-comparison}]
\begin{enumerate}
\item Fix an admissible public tuple and its matched smoothness tuple:
\[
v=(p,\alpha,\beta,\gamma)\in\mathcal V,
\qquad
w=(\alpha,\beta,\gamma),
\qquad
v_2=(2,\alpha,\beta,\gamma)\in\mathcal V.
\]
Use the quantities of \cref{obj:def:sharp-frontier-scales}, in particular
\[
q=\frac{p-1}{p},
\qquad
S=\alpha+\beta,
\qquad
E(p)=\min\{E_0(p),E_4(p)\}.
\]
Admissibility gives
\[
p-1>0,\qquad
0<q\le\frac12,\qquad
S>0,\qquad
\gamma>0,\qquad
2\gamma+q>0,\qquad
D_p>0.
\]
Indeed, \(q\le1/2\) follows by multiplying by \(p>0\) and using \(p\le2\); every term added to \(1\) in \(D_p\) is positive. Consequently,
\[
E_0(p)>0,\qquad E_4(p)>0,\qquad E(p)>0.
\]
The same conclusions apply at \(v_2\).

We first establish the phase selection needed for the lower bounds. Since
\[
pq=p-1,
\qquad
\frac1q-1=\frac1{p-1},
\]
subtracting the two channel exponents gives
\[
\begin{aligned}
E_4(p)-E_0(p)
&=\frac{2S(2\gamma+q)-2\gamma qD_p}
        {(2\gamma+q)D_p}\\
&=\frac{2\gamma q}{(2\gamma+q)D_p}
  \left(
  \frac{2S}{q}-2\alpha-\frac{\beta}{q}
  +\frac{S}{2\gamma}-1
  \right)\\
&=\frac{2\gamma q}{(2\gamma+q)D_p}\bigl(F(p)-1\bigr).
\end{aligned}
\]
The coefficient multiplying \(F(p)-1\) is strictly positive. Thus
\[
\begin{aligned}
F(p)\ge1&\ \Longrightarrow\ E_0(p)\le E_4(p)
                       \ \Longrightarrow\ E(p)=E_0(p),\\
F(p)<1&\ \Longrightarrow\ E_4(p)<E_0(p)
                       \ \Longrightarrow\ E(p)=E_4(p).
\end{aligned}
\]
These implications also hold with \(p=2\).

The definition of the lower multiplier gives
\[
N_{\mathrm{low}}(v)\ge2,
\qquad
c_0^{\mathrm e}>0,
\qquad
k_v>0,
\qquad
c_v>0,
\]
because every power occurring in either branch has a positive base. Applying this argument at \(v_2\), and applying \cref{obj:thm:original-record-attainable-rate} at both tuples, yields
\[
c_w^{\mathrm b}>0,\qquad C_v>0,\qquad C_w^{\mathrm b}>0.
\]
For every \(n\ge2\), positivity of real powers also gives
\[
\rho_n(v)>0,
\qquad
\rho_n^{\mathrm b}(w)>0.
\]
% lean: CausalSmith.Stat.FinitepHomogeneityDensegamma.phase_channel_difference
% lean: CausalSmith.Stat.FinitepHomogeneityDensegamma.phase_channel_selection
% lean: CausalSmith.Stat.FinitepHomogeneityDensegamma.cLower_pos
% lean: CausalSmith.Stat.FinitepHomogeneityDensegamma.matched_scales_pos

\item Fix an integer \(n\ge2\). We obtain the bounded lower bound by splitting according to the phase at \(v_2\).

If \(F(2)<1\), then
\[
E(2)=E_4(2),
\qquad
c_w^{\mathrm b}\rho_n^{\mathrm b}(w)
=c_w^{\mathrm b}n^{-E_4(2)}.
\]
The bounded conclusion of \cref{obj:lem:rough-full-record-sharp-lower} supplies normalized finite alternative witnesses and
\[
B_n^{\mathrm b}
 \bigl(w,c_w^{\mathrm b}\rho_n^{\mathrm b}(w)\bigr)
\ge\frac25.
\]
A normalized finite prior has a positive-weight atom: otherwise every nonnegative weight would be zero, contradicting that their sum is one. Choosing such an alternative atom gives
\[
\exists P\in\mathcal M_w^{\mathrm b}:
\quad
c_w^{\mathrm b}\rho_n^{\mathrm b}(w)\le d(P).
\]

If \(F(2)\ge1\), the public definitions and the phase selection give
\[
E(2)=E_0(2)=\frac{2\gamma}{4\gamma+1},
\qquad
c_w^{\mathrm b}=c_0^{\mathrm e}
=\frac{4^{-\gamma}}{16\sqrt3}.
\]
The signed-binary witnesses in \cref{obj:lem:paired-tent-full-record-lower} therefore give the same risk bound and a positive-weight alternative satisfying
\[
P\in\mathcal M_w^{\mathrm{bin}}
\subseteq\mathcal M_w^{\mathrm b},
\qquad
c_w^{\mathrm b}\rho_n^{\mathrm b}(w)\le d(P),
\]
where the inclusion follows from \cref{obj:prop:bounded-submodel-comparison}. Hence, in both cases,
\[
\exists P\in\mathcal M_w^{\mathrm b}:
\quad c_w^{\mathrm b}\rho_n^{\mathrm b}(w)\le d(P),
\qquad
B_n^{\mathrm b}
 \bigl(w,c_w^{\mathrm b}\rho_n^{\mathrm b}(w)\bigr)
\ge\frac25.
\]

We also verify that this bounded lower separation lies below the bounded cap. For any admissible tuple, the branch \(F(p)\ge1\) gives \(c_v=c_0^{\mathrm e}\). In the other branch, the phase identity gives \(E_0(p)-E_4(p)>0\), so
\[
N_{\mathrm{low}}(v)^{-(E_0(p)-E_4(p))}\le1,
\qquad
c_v\le
c_0^{\mathrm e}N_{\mathrm{low}}(v)^{-(E_0(p)-E_4(p))}
\le c_0^{\mathrm e}.
\]
Since \(E(p)>0\) and \(n\ge2\),
\[
0<\rho_n(v)\le1,
\qquad
0<c_v\rho_n(v)\le c_0^{\mathrm e}
\le\frac1{16\sqrt3}
<\frac1{8\sqrt2}=d_0.
\]
The strict numerical inequality follows from
\(\sqrt2<2\sqrt3\), obtained by squaring these positive quantities. Applying this bound at \(v_2\), and using
\(d_0\le D_w^{\mathrm b}\) from \cref{obj:prop:bounded-submodel-comparison}, gives
\[
0<c_w^{\mathrm b}\rho_n^{\mathrm b}(w)
<d_0\le D_w^{\mathrm b}.
\]

For the original model, \cref{obj:thm:tail-essential-full-record-converse} supplies
\[
0<c_v\rho_n(v)<D_v,
\qquad
B_n\bigl(v,c_v\rho_n(v)\bigr)\ge\frac25,
\]
and a normalized finite alternative prior. Choosing a positive-weight atom as above gives
\[
\exists P\in\mathcal M_v:
\quad c_v\rho_n(v)\le d(P)<D_v.
\]

We now pass from these failures to the capped radii. For
\(0<r_1\le r_2<D_v\), with a nonempty alternative class at \(r_2\), inclusion of the alternative classes gives, for every admissible test,
\[
\sup_{\substack{P\in\mathcal M_v\\ d(P)\ge r_2}}
 \{1-\mathsf R_n(P,\phi)\}
\le
\sup_{\substack{P\in\mathcal M_v\\ d(P)\ge r_1}}
 \{1-\mathsf R_n(P,\phi)\}.
\]
The errors lie in \([0,1]\), and the null-size constraint is unchanged. Taking infima over tests therefore yields
\[
B_n(v,r_2)\le B_n(v,r_1).
\]
The identical argument gives bounded-risk monotonicity.

If a successful separation \(\widetilde r\in(0,D_v)\) satisfied
\(\widetilde r<c_v\rho_n(v)\), monotonicity and its success condition would imply
\[
\frac25
\le B_n\bigl(v,c_v\rho_n(v)\bigr)
\le B_n(v,\widetilde r)
\le\frac1{10},
\]
a contradiction. The sentinel \(D_v\) also exceeds \(c_v\rho_n(v)\). Thus every element of the nonempty set defining the capped infimum in \cref{obj:def:critical-radius} is at least \(c_v\rho_n(v)\). The bounded argument uses its positive-weight witness and the legal separation established above. Consequently,
\[
c_v\rho_n(v)\le r_n^*(v),
\qquad
c_w^{\mathrm b}\rho_n^{\mathrm b}(w)
\le r_n^{*,\mathrm b}(w).
\]
The displayed risk bounds also imply both asserted strict inequalities against \(1/10\).
% lean: CausalSmith.Stat.FinitepHomogeneityDensegamma.bounded_matched_lower
% lean: CausalSmith.Stat.FinitepHomogeneityDensegamma.matched_lower_separation
% lean: CausalSmith.Stat.FinitepHomogeneityDensegamma.normalized_prior_positive_atom
% lean: CausalSmith.Stat.FinitepHomogeneityDensegamma.cappedRadius_lower_of_nonempty_failure

\item By \cref{obj:thm:original-record-attainable-rate}, the public tests satisfy
\[
r_n^*(v)\le C_v\rho_n(v),
\qquad
r_n^{*,\mathrm b}(w)\le C_w^{\mathrm b}\rho_n^{\mathrm b}(w).
\]
Together with \cref{obj:prop:bounded-submodel-comparison} and the preceding lower bounds, this gives
\[
\begin{aligned}
c_v\rho_n(v)
&\le r_n^*(v)\le C_v\rho_n(v),\\
c_w^{\mathrm b}\rho_n^{\mathrm b}(w)
&\le r_n^{*,\mathrm b}(w)
=r_n^{*,\mathrm{bin}}(w)
\le C_w^{\mathrm b}\rho_n^{\mathrm b}(w).
\end{aligned}
\]
The same attainment result supplies
\[
\begin{aligned}
\mathsf R_n(P,\phi^*_{n,v})&\le0.0075<\frac1{10}
&&\quad(P\in H_0(v)),\\
\mathsf R_n(P,\phi^{*,\mathrm b}_{n,w})&\le0.0075<\frac1{10}
&&\quad(P\in H_0^{\mathrm b}(w)).
\end{aligned}
\]
Under the respective conditions
\[
C_v\rho_n(v)<D_v,
\qquad
C_w^{\mathrm b}\rho_n^{\mathrm b}(w)<D_w^{\mathrm b},
\]
it supplies, respectively,
\[
\begin{aligned}
1-\mathsf R_n(P,\phi^*_{n,v})&\le0.0075<\frac1{10}
&&\quad(P\in\mathcal M_v,\ d(P)\ge C_v\rho_n(v)),\\
1-\mathsf R_n(P,\phi^{*,\mathrm b}_{n,w})&\le0.0075<\frac1{10}
&&\quad(P\in\mathcal M_w^{\mathrm b},\
d(P)\ge C_w^{\mathrm b}\rho_n^{\mathrm b}(w)).
\end{aligned}
\]

The bounded radius is positive by its lower bound. Dividing the matched bounds therefore gives
\[
\frac{c_v\rho_n(v)}
     {C_w^{\mathrm b}\rho_n^{\mathrm b}(w)}
\le\Delta_n(v)
\le
\frac{C_v\rho_n(v)}
     {c_w^{\mathrm b}\rho_n^{\mathrm b}(w)}.
\]
For \(n>0\), the law of real powers gives
\[
\frac{\rho_n(v)}{\rho_n^{\mathrm b}(w)}
=\frac{n^{-E(p)}}{n^{-E(2)}}
=n^{E(2)-E(p)}.
\]
Hence
\[
\frac{c_v}{C_w^{\mathrm b}}n^{E(2)-E(p)}
\le\Delta_n(v)
\le
\frac{C_v}{c_w^{\mathrm b}}n^{E(2)-E(p)}.
\]
Finally, \cref{obj:prop:bounded-submodel-comparison} gives
\[
0<r_n^{*,\mathrm b}(w)\le r_n^*(v),
\qquad
1\le\Delta_n(v).
\]
% lean: CausalSmith.Stat.FinitepHomogeneityDensegamma.matched_radius_ratio_bounds
% lean: CausalSmith.Stat.FinitepHomogeneityDensegamma.frontier_at

\item Since \(E(p)>0\) and \(E(2)>0\), the definitions of the scales give
\[
\rho_n(v)=n^{-E(p)}\longrightarrow0,
\qquad
\rho_n^{\mathrm b}(w)=n^{-E(2)}\longrightarrow0.
\]
The public thresholds are those of \cref{obj:def:sharp-frontier-scales}, evaluated at \(v\) and \(v_2\). Their formulas in \cref{obj:thm:original-record-attainable-rate} give
\[
N_v\ge4,\qquad N_w^{\mathrm b}\ge4.
\]
Its witness-threshold conclusion therefore applies to every natural number satisfying the corresponding threshold and yields
\[
\begin{aligned}
n\ge N_v&\ \Longrightarrow\ C_v\rho_n(v)<d_0,\\
n\ge N_w^{\mathrm b}&\ \Longrightarrow\
C_w^{\mathrm b}\rho_n^{\mathrm b}(w)<d_0.
\end{aligned}
\]
% lean: CausalSmith.Stat.FinitepHomogeneityDensegamma.rho_tendsto_zero
% lean: CausalSmith.Stat.FinitepHomogeneityDensegamma.rhoBounded_tendsto_zero
% lean: CausalSmith.Stat.FinitepHomogeneityDensegamma.frontier_at

\item We make the exponent difference explicit, always preserving \(w\) when evaluating at \(2\). Define
\[
D_2=1+2S+\frac{S}{2\gamma},
\qquad
F_2=2S+\frac{S}{2\gamma}.
\]
Substitution of \(q=(p-1)/p\) gives
\[
D_p=D_2+\frac{\beta(2-p)}{p-1},
\qquad
F(p)=F_2+\frac{(2\alpha+\beta)(2-p)}{p-1}
\ge F_2.
\]
For the supplied-propensity channel, a common denominator yields
\[
\begin{aligned}
E_0(2)-E_0(p)
&=\frac{2\gamma}{4\gamma+1}
 -\frac{2\gamma(p-1)}{2\gamma p+p-1}\\
&=\frac{4\gamma^2(2-p)}
 {(4\gamma+1)(2\gamma p+p-1)}.
\end{aligned}
\]
Both denominators are positive. For the rough channel, the identity for \(D_p\) gives
\[
\begin{aligned}
E_4(2)-E_4(p)
&=\frac{2S}{D_2}-\frac{2S}{D_p}\\
&=\frac{2S(D_p-D_2)}{D_2D_p}\\
&=\frac{2S\beta(2-p)}{(p-1)D_pD_2}.
\end{aligned}
\]
Combining these identities with the phase rules \(E(p)=E_0(p)\) for \(F(p)\ge1\) and \(E(p)=E_4(p)\) for \(F(p)<1\), applied also at \(2\), gives the exhaustive expression
\[
E(2)-E(p)=
\begin{cases}
\displaystyle
\frac{4\gamma^2(2-p)}
 {(4\gamma+1)(2\gamma p+p-1)},
&F_2\ge1,\\[6pt]
\displaystyle
E_4(2)-E_0(p),
&F_2<1\le F(p),\\[4pt]
\displaystyle
\frac{2S\beta(2-p)}{(p-1)D_pD_2},
&F(p)<1.
\end{cases}
\]
Indeed, in the first case \(F(p)\ge F_2\ge1\), so both minimum exponents select the supplied-propensity channel. In the second case the bounded exponent selects the rough channel and the finite-moment exponent selects the supplied-propensity channel. In the third case \(F_2\le F(p)<1\), so both select the rough channel.

If \(p<2\), the two channel differences displayed above are strictly positive. Therefore
\[
E(p)\le E_0(p)<E_0(2),
\qquad
E(p)\le E_4(p)<E_4(2),
\]
and hence
\[
E(p)<\min\{E_0(2),E_4(2)\}=E(2).
\]
Using the scale definitions in \cref{obj:def:sharp-frontier-scales}, we obtain
\[
\frac{\rho_n(v)}{\rho_n^{\mathrm b}(w)}
=n^{E(2)-E(p)}
\longrightarrow\infty.
\]
% lean: CausalSmith.Stat.FinitepHomogeneityDensegamma.phase_bounded_identities
% lean: CausalSmith.Stat.FinitepHomogeneityDensegamma.phase_bounded_tail_difference
% lean: CausalSmith.Stat.FinitepHomogeneityDensegamma.phase_bounded_rough_difference
% lean: CausalSmith.Stat.FinitepHomogeneityDensegamma.phase_difference_eq
% lean: CausalSmith.Stat.FinitepHomogeneityDensegamma.exponent_phase_algebra

\item Suppose \(p<2\). Apply \cref{obj:thm:tail-essential-full-record-converse} to the selected priors and magnitudes of \cref{obj:def:converse-handle}. It gives
\[
L_{n,v}\longrightarrow\infty.
\]
For every \(n\ge2\), both selected priors are normalized finite probability priors, and their positive-weight laws satisfy
\[
\begin{aligned}
\pi_{0,n,v}(\{P\})>0
&\ \Longrightarrow\ P\in H_0(v),\\
\pi_{1,n,v}(\{P\})>0
&\ \Longrightarrow\
P\in\mathcal M_v,\quad
c_v\rho_n(v)\le d(P)<D_v.
\end{aligned}
\]
For each positive-weight law in either prior, their conditional outcome kernels satisfy
\[
Q_{a,P}\bigl(\{0,-L_{n,v},L_{n,v}\}\mid x\bigr)=1
\]
for \(a\in\{0,1\}\) and \(\lambda\)-almost every \(x\). The same result supplies the complete-record bound
\[
\operatorname{TV}(\mathbb P_{0,n,v},\mathbb P_{1,n,v})
<\frac12.
\]

For any normalized finite probability prior \(\pi\), write its complete-record mixture as
\[
\mathbb P_{\pi,n}=\int P^{\otimes n}\,\pi(dP).
\]
The bounded-prior conclusion of the cited result gives a threshold \(N(v)\in\mathbb N\) such that, for every \(n\ge N(v)\) and every pair of normalized finite probability priors \(\pi_0,\pi_1\),
\[
\left.
\begin{aligned}
\pi_0(\{P\})>0&\ \Longrightarrow\ P\in H_0^{\mathrm b}(w),\\
\pi_1(\{P\})>0&\ \Longrightarrow\
P\in\mathcal M_w^{\mathrm b},\quad c_v\rho_n(v)\le d(P)
\end{aligned}
\right\}
\ \Longrightarrow\
\operatorname{TV}(\mathbb P_{\pi_0,n},\mathbb P_{\pi_1,n})
\ge\frac12.
\]
These are precisely the asserted support, magnitude, and complete-record mixture conclusions.
% lean: CausalSmith.Stat.FinitepHomogeneityDensegamma.frontier_at

\item For \(p<2\), the lower ratio bound and the strictly positive exponent difference give
\[
0<\frac{c_v}{C_w^{\mathrm b}},
\qquad
\frac{c_v}{C_w^{\mathrm b}}n^{E(2)-E(p)}
\longrightarrow\infty,
\qquad
\Delta_n(v)\longrightarrow\infty.
\]
The definition in \cref{obj:def:tail-region} includes admissibility and \(p<2\). The preceding divergence supplies its remaining condition, so
\[
\mathcal T=\{v\in\mathcal V:p<2\}.
\]

To characterize \cref{obj:def:equality-region}, suppose that an admissible tuple has a bounded radius ratio. Thus there exists \(C_\Delta\in\mathbb R\) such that
\[
\Delta_n(v)\le C_\Delta
\qquad(n\ge2).
\]
If \(p\ne2\), admissibility implies \(p<2\), and the divergence just established supplies \(N_\Delta\in\mathbb N\) with
\[
n\ge N_\Delta\ \Longrightarrow\ \Delta_n(v)>C_\Delta.
\]
At
\[
n_\Delta=\max\{2,N_\Delta\},
\]
the two bounds would give
\[
C_\Delta<\Delta_{n_\Delta}(v)\le C_\Delta,
\]
a contradiction. Thus boundedness implies \(p=2\). Conversely, when \(p=2\), evaluation at \(2\) preserves the entire tuple, and the upper ratio bound reduces to
\[
\Delta_n(v)\le\frac{C_v}{c_w^{\mathrm b}}
\qquad(n\ge2).
\]
Consequently,
\[
\mathcal E=\{v\in\mathcal V:p=2\}.
\]
% lean: CausalSmith.Stat.FinitepHomogeneityDensegamma.radiusRatio_tendsto_atTop
% lean: CausalSmith.Stat.FinitepHomogeneityDensegamma.tailRegion_characterization
% lean: CausalSmith.Stat.FinitepHomogeneityDensegamma.second_moment_radiusRatio_bounded
% lean: CausalSmith.Stat.FinitepHomogeneityDensegamma.equalityRegion_characterization

\item Define the ambient set
\[
\mathcal U
=\{(p,\alpha,\beta,\gamma)\in\mathbb R^4:p<2\}.
\]
The coordinate map to \(p\) is continuous, so \(\mathcal U\) is open in the product topology. The tuple
\[
v_{\mathrm o}=(3/2,1/2,1/2,1/2)
\]
satisfies every admissibility inequality and belongs to \(\mathcal U\). Thus
\[
v_{\mathrm o}\in\mathcal U\cap\mathcal V,
\qquad
\mathcal U\cap\mathcal V\ne\varnothing.
\]
The characterization of \(\mathcal T\) gives
\[
\mathcal U\cap\mathcal V\subseteq\mathcal T,
\]
which supplies the required open set.
% lean: CausalSmith.Stat.FinitepHomogeneityDensegamma.tailRegion_nonempty_open

\item It remains to prove compact uniformity. For each admissible tuple define the auxiliary positive base and lower envelope by
\[
\Theta(v)
=\max\left\{2,32^{D_p\gamma/(2S)}\right\}+1,
\]
\[
\underline c_{\mathrm{env}}(v)
=\min\left\{
c_0^{\mathrm e},\,
k_v,\,
c_0^{\mathrm e}\Theta(v)^{-|E_0(p)-E_4(p)|}
\right\}.
\]
Every entry in this minimum is positive, so
\[
\underline c_{\mathrm{env}}(v)>0.
\]
The ceiling bound gives
\[
2\le N_{\mathrm{low}}(v)\le\Theta(v).
\]
Since the exponent \(-|E_0(p)-E_4(p)|\) is nonpositive, and since
\(N_{\mathrm{low}}(v)\ge1\), monotonicity in the base and then in the exponent gives
\[
\begin{aligned}
\Theta(v)^{-|E_0(p)-E_4(p)|}
&\le N_{\mathrm{low}}(v)^{-|E_0(p)-E_4(p)|}\\
&\le N_{\mathrm{low}}(v)^{-(E_0(p)-E_4(p))}.
\end{aligned}
\]
On \(F(p)\ge1\), the envelope is at most
\(c_0^{\mathrm e}=c_v\). On \(F(p)<1\), it is at most each of
\[
k_v,
\qquad
c_0^{\mathrm e}
N_{\mathrm{low}}(v)^{-(E_0(p)-E_4(p))}.
\]
The branchwise definition therefore gives, throughout the admissible domain,
\[
0<\underline c_{\mathrm{env}}(v)\le c_v.
\]

The envelope is continuous at every admissible tuple. Indeed, the coordinate maps are continuous; the denominators defining \(q,D_p,E_0,E_4\) are nonzero there; and every variable power in the envelope has a positive base. Absolute value, maximum, minimum, and the remaining arithmetic operations preserve continuity.

Let \(\mathcal K\subseteq\mathcal V\) be compact and nonempty. The continuous envelope attains a minimum at some \(v_{\mathcal K}\in\mathcal K\). Define
\[
c_{\mathcal K}
=\underline c_{\mathrm{env}}(v_{\mathcal K})
=\min_{v\in\mathcal K}\underline c_{\mathrm{env}}(v)>0.
\]
Then
\[
c_{\mathcal K}
\le\underline c_{\mathrm{env}}(v)
\le c_v
\qquad(v\in\mathcal K).
\]
The compact-uniformity conclusion of \cref{obj:thm:original-record-attainable-rate} supplies real constants \(C_{\mathcal K},N_{\mathcal K}\) such that
\[
C_v\le C_{\mathcal K},
\qquad
N_v\le N_{\mathcal K}
\qquad(v\in\mathcal K).
\]
For the empty compact set, the choices
\[
c_{\mathcal K}=1,\qquad C_{\mathcal K}=0,\qquad N_{\mathcal K}=0
\]
satisfy all requirements vacuously.

Finally, let \(\mathcal K_{\mathrm b}\subseteq\mathcal W\) be compact. Its embedded parameter set is
\[
\mathcal K_2
=\{(2,\alpha,\beta,\gamma):
(\alpha,\beta,\gamma)\in\mathcal K_{\mathrm b}\}.
\]
The embedding \(w\mapsto(2,w)\) is continuous, so \(\mathcal K_2\) is compact, and all its tuples are admissible. Applying the preceding result to \(\mathcal K_2\), and using the public identities
\[
c_w^{\mathrm b}=c_{(2,w)},
\qquad
C_w^{\mathrm b}=C_{(2,w)},
\qquad
N_w^{\mathrm b}=N_{(2,w)},
\]
gives
\[
c_{\mathcal K_2}\le c_w^{\mathrm b},
\qquad
C_w^{\mathrm b}\le C_{\mathcal K_2},
\qquad
N_w^{\mathrm b}\le N_{\mathcal K_2}
\qquad(w\in\mathcal K_{\mathrm b}),
\]
with \(c_{\mathcal K_2}>0\). This establishes both compact-uniformity assertions.
% lean: CausalSmith.Stat.FinitepHomogeneityDensegamma.compactLowerEnvelope_le
% lean: CausalSmith.Stat.FinitepHomogeneityDensegamma.continuousAt_compactLowerEnvelope
% lean: CausalSmith.Stat.FinitepHomogeneityDensegamma.cLower_compact_uniformity
% lean: CausalSmith.Stat.FinitepHomogeneityDensegamma.frontier_compact_uniformity
\end{enumerate}
\end{proof}

\begin{proof}[Proof of \cref{obj:thm:oracle-frontier-broad-class}]
Fix an admissible tuple \(v\). We use the constants and exponent defined in the statement.

\begin{enumerate}
\item Comparing the restrictions in
\cref{obj:def:model,obj:def:oracle-broad-model,obj:def:null,obj:def:oracle-broad-null}
gives
\[
\mathcal M_v\subseteq\mathcal M_v^{\mathrm{e,br}},
\qquad
H_0(v)\subseteq H_0^{\mathrm{e,br}}(v).
\]
The law displayed in \cref{obj:lem:causal-nonempty} therefore supplies a member of the broader model with distance \(d_0\).

For any \(P\in\mathcal M_v^{\mathrm{e,br}}\), put
\[
\bar\tau_P=\int_{[0,1]}\tau_P(x)\,\lambda(dx).
\]
Expanding the centered square and using the effect cap gives
\[
d(P)^2
=\int_{[0,1]}\tau_P(x)^2\,\lambda(dx)-\bar\tau_P^2
\le\int_{[0,1]}\tau_P(x)^2\,\lambda(dx)
\le\frac14.
\]
Thus the set of broader-model distances is nonempty and bounded above, and
\[
d_0\le D_v^{\mathrm{e,br}}\le\frac12.
\]
Furthermore, \(\gamma\le1\) implies
\[
4^{-\gamma}\ge4^{-1}=\frac14,
\qquad
c_v^{\mathrm e}
=\frac{4^{-\gamma}}{16\sqrt3}
\ge\frac1{64\sqrt3},
\qquad
C_v^{\mathrm e}>0.
\]
% lean: CausalSmith.Stat.FinitepHomogeneityDensegamma.oracle_model_distance_bounds
% lean: CausalSmith.Stat.FinitepHomogeneityDensegamma.cOracle_uniform_lower
% lean: CausalSmith.Stat.FinitepHomogeneityDensegamma.COr_pos

\item Fix an integer \(n\ge2\), and define
\[
r_{\mathrm{low}}=c_v^{\mathrm e}\rho_n^{\mathrm e}(v).
\]
By \cref{obj:lem:paired-tent-full-record-lower}, there are a null law \(P_0\in H_0(v)\) and a normalized finite prior
\[
\pi_1=\sum_{j=1}^k\omega_j\delta_{P_j},
\qquad
\omega_j\ge0,
\qquad
\sum_{j=1}^k\omega_j=1,
\]
with at least one positive weight, such that
\[
e_{P_0}(x)=\frac12,
\qquad
\omega_j>0
\ \Longrightarrow\
P_j\in\mathcal M_v,\quad
e_{P_j}(x)=\frac12,\quad
r_{\mathrm{low}}\le d(P_j)<d_0.
\]
Define the supplied function and the complete-record mixture by
\[
f=e_{P_0},
\qquad
\overline P_1^{(n)}
=\sum_{j=1}^k\omega_jP_j^{\otimes n}.
\]
The function \(f\) is continuous and satisfies \(f(x)=1/2\) everywhere. The same cited result gives
\[
\operatorname{TV}\!\left(P_0^{\otimes n},\overline P_1^{(n)}\right)<\frac12.
\]
The inclusions already established place \(P_0\) in the broader null and every positive-weight \(P_j\) in the broader model, with \(e_{P_j}=f\).

Admissibility gives \(q>0\), \(\gamma>0\), and hence \(E_0>0\). Consequently,
\[
0<r_{\mathrm{low}}
\le d(P_j)<d_0\le D_v^{\mathrm{e,br}}
\qquad(\omega_j>0).
\]
In particular, \(r_{\mathrm{low}}\) is a legal separation with a nonempty alternative class.
% lean: CausalSmith.Stat.FinitepHomogeneityDensegamma.paired_tent_full_record_lower
% lean: CausalSmith.Stat.FinitepHomogeneityDensegamma.oracle_frontier_broad_class

\item Let \(\phi^{\mathrm e}\) be an allowed supplied-propensity test satisfying the broader-null size constraint. Freeze its function input at \(f\), and average its public randomization:
\[
\overline\phi_f(\mathcal D_n)
=\int_0^1\phi^{\mathrm e}(f,\mathcal D_n,U)\,dU,
\qquad
0\le\overline\phi_f(\mathcal D_n)\le1.
\]
Joint measurability makes this a measurable function of the sample. Since all positive-weight witnesses have supplied function \(f\), integration against the finite mixture gives
\[
\int\overline\phi_f\,dP_0^{\otimes n}
=\mathsf R_n^{\mathrm e}(P_0,\phi^{\mathrm e}),
\qquad
\int\overline\phi_f\,d\overline P_1^{(n)}
=\sum_{j=1}^k\omega_j\mathsf R_n^{\mathrm e}(P_j,\phi^{\mathrm e}).
\]
For a measurable function taking values in \([0,1]\), integrating its level-set probability differences bounds the difference of its expectations by total variation. Therefore
\[
\left|
\sum_{j=1}^k\omega_j\mathsf R_n^{\mathrm e}(P_j,\phi^{\mathrm e})
-\mathsf R_n^{\mathrm e}(P_0,\phi^{\mathrm e})
\right|
\le
\operatorname{TV}\!\left(P_0^{\otimes n},\overline P_1^{(n)}\right),
\]
and hence
\[
\sum_{j=1}^k\omega_j\mathsf R_n^{\mathrm e}(P_j,\phi^{\mathrm e})
<\frac1{10}+\frac12=\frac35.
\]
If every positive-weight atom had type-II error strictly below \(2/5\), then its rejection expectation would exceed \(3/5\). Zero-weight atoms contribute zero, so normalization would imply
\[
\sum_{j=1}^k\omega_j\mathsf R_n^{\mathrm e}(P_j,\phi^{\mathrm e})
\ge\frac35\sum_{j=1}^k\omega_j
=\frac35,
\]
a contradiction. Thus
\[
\exists j:\quad
\omega_j>0,\qquad
1-\mathsf R_n^{\mathrm e}(P_j,\phi^{\mathrm e})\ge\frac25.
\]

The feasible test class is nonempty: the jointly measurable test
\[
\phi_0^{\mathrm e}(f',\mathcal D_n,U)=0
\qquad
\bigl(f'\in C([0,1]),\ \mathcal D_n\text{ an }n\text{-record sample},\
U\in[0,1]\bigr)
\]
has rejection expectation zero under every law. For every allowed test and every law,
\[
0\le\mathsf R_n^{\mathrm e}(P,\phi^{\mathrm e})\le1.
\]
Thus the alternative errors are bounded above by one. The positive-weight atom just obtained belongs to the alternative class at \(r_{\mathrm{low}}\), so
\[
\sup_{\substack{P\in\mathcal M_v^{\mathrm{e,br}}\\
d(P)\ge r_{\mathrm{low}}}}
\left\{1-\mathsf R_n^{\mathrm e}(P,\phi^{\mathrm e})\right\}
\ge\frac25.
\]
Taking the infimum over feasible tests in
\cref{obj:def:oracle-broad-risk} proves
\[
B_n^{\mathrm{e,br}}(v,r_{\mathrm{low}})\ge\frac25.
\]
This also proves the required error witness for the particular prior constructed above.
% lean: CausalSmith.Stat.FinitepHomogeneityDensegamma.oracle_single_null_prior_error_witness
% lean: CausalSmith.Stat.FinitepHomogeneityDensegamma.oracleRejectProb_bounds
% lean: CausalSmith.Stat.FinitepHomogeneityDensegamma.oracle_frontier_broad_class

\item For every separation \(r\) with \(0<r\le r_{\mathrm{low}}\), the alternative class at \(r_{\mathrm{low}}\) is contained in the alternative class at \(r\). Both are nonempty, their error suprema are bounded, and the feasible test class is unchanged. Therefore
\[
B_n^{\mathrm{e,br}}(v,r)
\ge B_n^{\mathrm{e,br}}(v,r_{\mathrm{low}})
\ge\frac25>\frac1{10}.
\]
Every successful separation in the set defining
\cref{obj:def:oracle-broad-radius} is consequently above \(r_{\mathrm{low}}\); its sentinel \(D_v^{\mathrm{e,br}}\) is also above \(r_{\mathrm{low}}\). Taking the infimum gives
\[
r_n^{*,\mathrm{e,br}}(v)\ge r_{\mathrm{low}}.
\]
% lean: CausalSmith.Stat.FinitepHomogeneityDensegamma.oracleBroadCriticalRadius_ge_of_failed

\item We now establish calibration and power directly on the broader model. Use the construction in \cref{obj:def:oracle-handle}, with
\[
s=\lfloor n/2\rfloor,\qquad
a=s^{-E_0},\qquad
T=2^{\lceil\log_2(a^{-1/(p-1)})\rceil},\qquad
J=2^{\lceil\log_2(a^{-1/\gamma})\rceil},
\]
\[
b=20T^{1-p},
\qquad
\Lambda=\frac{80T^{2-p}}s.
\]
Here
\[
s\ge1,\qquad 0<a\le1,\qquad T\ge1,\qquad J\ge1,\qquad
b\ge0,\qquad\Lambda>0.
\]
For a broader-model law \(P\), overlap makes the clipped supplied propensity equal to \(e_P\).

The elementary clipping inequalities are
\[
|y-\ell_T(y)|\le |y|^pT^{1-p},
\qquad
\ell_T(y)^2\le |y|^pT^{2-p}.
\]
Indeed, when \(|y|\le T\), the first difference vanishes and
\[
y^2=|y|^p|y|^{2-p}\le |y|^pT^{2-p}.
\]
When \(|y|>T\), the clipping remainder is at most \(|y|\), and
\[
|y|
=|y|^p|y|^{1-p}
\le |y|^pT^{1-p},
\qquad
\ell_T(y)^2\le T^2
=T^pT^{2-p}\le |y|^pT^{2-p}.
\]
The raw moment restriction gives first-moment integrability in both arms, using \(|y|\le1+|y|^p\). Integrating the two clipping inequalities yields, for \(\lambda\)-almost every \(x\),
\[
\left|\int_{\mathbb R}(y-\ell_T(y))\,Q_{\iota_{\mathrm{arm}},P}(dy\mid x)\right|
\le10T^{1-p},
\qquad
\int_{\mathbb R}\ell_T(y)^2\,Q_{\iota_{\mathrm{arm}},P}(dy\mid x)
\le10T^{2-p}
\quad(\iota_{\mathrm{arm}}\in\{0,1\}).
\]

Define the clipped conditional score mean by
\[
g_{T,P}(x)
=
\int_{\mathbb R}\ell_T(y)\,Q_{1,P}(dy\mid x)
-
\int_{\mathbb R}\ell_T(y)\,Q_{0,P}(dy\mid x).
\]
Cancellation of the arm probabilities against the inverse-propensity denominators identifies this function as the conditional mean of \(R_i\). The conditional arm-mean identities then give
\[
|g_{T,P}(x)-\tau_P(x)|\le20T^{1-p}=b
\]
almost everywhere. The conditional second moment of the score satisfies
\[
\begin{aligned}
\mathbb E(R_i^2\mid X_i=x)
&=
\frac{\int\ell_T(y)^2\,Q_{1,P}(dy\mid x)}{e_P(x)}
+
\frac{\int\ell_T(y)^2\,Q_{0,P}(dy\mid x)}{1-e_P(x)}
\\
&\le40T^{2-p}+40T^{2-p}
=80T^{2-p}.
\end{aligned}
\]
% lean: CausalSmith.Stat.FinitepHomogeneityDensegamma.oracle_clipping_moments
% lean: CausalSmith.Stat.FinitepHomogeneityDensegamma.conditional_truncation_moments

\item Define the constant-removing matrix and the sample expectation notation by
\[
u=J^{-1/2}(1,\ldots,1)^T,\qquad
C_J=\operatorname{Id}_{\mathbb R^J}-uu^T,
\]
\[
\mathbb P_{n,P}=P^{\otimes n},
\qquad
\mathbb E_{n,P}[\cdot]=\int(\cdot)\,dP^{\otimes n}.
\]
The histogram cells in \cref{obj:def:projection} are disjoint and have Lebesgue mass \(1/J\). Consequently,
\[
\int_{[0,1]}F_J(x)F_J(x)^T\,\lambda(dx)
=\operatorname{Id}_{\mathbb R^J},
\qquad
u^TF_J(x)=1,
\qquad
C_JF_J(x)=F_J(x)-u.
\]
For every \(z_{\mathrm{dir}}\in\mathbb R^J\),
\[
\|C_Jz_{\mathrm{dir}}\|^2
=\|z_{\mathrm{dir}}\|^2-(u^Tz_{\mathrm{dir}})^2
\le\|z_{\mathrm{dir}}\|^2.
\]
Combining these identities with the conditional score second-moment bound gives
\[
\begin{aligned}
\mathbb E_{n,P}
\!\left[\bigl(z_{\mathrm{dir}}^T(F_J(X_i)-u)R_i\bigr)^2\right]
&\le80T^{2-p}
\int_{[0,1]}
\bigl(z_{\mathrm{dir}}^TC_JF_J(x)\bigr)^2\,\lambda(dx)
\\
&=80T^{2-p}\|C_Jz_{\mathrm{dir}}\|^2
\\
&\le80T^{2-p}\|z_{\mathrm{dir}}\|^2.
\end{aligned}
\]

The two disjoint evaluation blocks have equal size and consist of independent records. Their common mean and centered vectors are
\[
\mu_P
=\mathbb E_{n,P}V_1
=\mathbb E_{n,P}V_2
=C_J\int_{[0,1]}F_J(x)g_{T,P}(x)\,\lambda(dx),
\qquad
\mathsf Z_{\iota_{\mathrm{block}},P}=V_{\iota_{\mathrm{block}}}-\mu_P\quad(\iota_{\mathrm{block}}\in\{1,2\}).
\]
Variance additivity within a block and the preceding single-record bound imply
\[
\mathbb E_{n,P}
\bigl[(z_{\mathrm{dir}}^T\mathsf Z_{\iota_{\mathrm{block}},P})^2\bigr]
=
\operatorname{Var}_{n,P}(z_{\mathrm{dir}}^TV_{\iota_{\mathrm{block}}})
\le\Lambda\|z_{\mathrm{dir}}\|^2.
\]
The centered vectors are independent and have mean zero.

Define the quadratic statistic by
\[
W_P=\langle V_1,V_2\rangle.
\]
Block independence gives
\[
\mathbb E_{n,P}W_P=\|\mu_P\|^2,
\qquad
W_P-\|\mu_P\|^2
=
\langle\mu_P,\mathsf Z_{1,P}\rangle
+\langle\mu_P,\mathsf Z_{2,P}\rangle
+\langle\mathsf Z_{1,P},\mathsf Z_{2,P}\rangle.
\]
All three cross-product expectations in the square of this expansion vanish:
\[
\begin{aligned}
\mathbb E_{n,P}
\bigl[\langle\mu_P,\mathsf Z_{1,P}\rangle
      \langle\mu_P,\mathsf Z_{2,P}\rangle\bigr]&=0,\\
\mathbb E_{n,P}
\bigl[\langle\mu_P,\mathsf Z_{1,P}\rangle
      \langle\mathsf Z_{1,P},\mathsf Z_{2,P}\rangle\bigr]&=0,\\
\mathbb E_{n,P}
\bigl[\langle\mu_P,\mathsf Z_{2,P}\rangle
      \langle\mathsf Z_{1,P},\mathsf Z_{2,P}\rangle\bigr]&=0.
\end{aligned}
\]
The first equality uses independence and both zero means; the other two follow by integrating first over the block appearing linearly.

For the remaining quadratic term, independence and the directional second-moment bound give
\[
\begin{aligned}
\mathbb E_{n,P}
\bigl[\langle\mathsf Z_{1,P},\mathsf Z_{2,P}\rangle^2\bigr]
&\le\Lambda\,\mathbb E_{n,P}\|\mathsf Z_{1,P}\|^2\\
&=\Lambda\sum_{j=1}^J
\mathbb E_{n,P}(\mathsf Z_{1,P})_j^2
\le J\Lambda^2.
\end{aligned}
\]
Thus
\[
\operatorname{Var}_{n,P}(W_P)
\le2\Lambda\|\mu_P\|^2+J\Lambda^2.
\]
All quantities have finite second moments, since the clipped scores and finite histogram features are bounded.

Define the good event by
\[
\mathcal G_{n,P}
=
\left\{\mathcal D_n:
|W_P-\|\mu_P\|^2|
\le16\bigl(\sqrt\Lambda\,\|\mu_P\|+\sqrt J\,\Lambda\bigr)
\right\}.
\]
Its tolerance is strictly positive. Applying the second-moment tail bound gives
\[
\begin{aligned}
\mathbb P_{n,P}(\mathcal G_{n,P}^{\,c})
&\le
\frac{2\Lambda\|\mu_P\|^2+J\Lambda^2}
{256\bigl(\sqrt\Lambda\,\|\mu_P\|+\sqrt J\,\Lambda\bigr)^2}\\
&\le\frac1{128}<0.01.
\end{aligned}
\]
The last inequality uses the nonnegative cross term in the denominator.
% lean: CausalSmith.Stat.FinitepHomogeneityDensegamma.oracleBlockVec_variance_le
% lean: CausalSmith.Stat.FinitepHomogeneityDensegamma.oracleBlockVec_inner_integral
% lean: CausalSmith.Stat.FinitepHomogeneityDensegamma.oracleBlockVec_quadratic_variance_le
% lean: CausalSmith.Stat.FinitepHomogeneityDensegamma.oracleBlockVec_tail_le

\item Suppose \(P\in H_0^{\mathrm{e,br}}(v)\), and choose its null constant:
\[
c_{\mathrm{null},P}\in[-1/2,1/2],
\qquad
\tau_P(x)=c_{\mathrm{null},P}\quad(x\in[0,1]).
\]
The histogram coefficient of this constant is \(c_{\mathrm{null},P}u\), which \(C_J\) annihilates. Hence
\[
\mu_P
=C_J\int_{[0,1]}
F_J(x)\bigl(g_{T,P}(x)-c_{\mathrm{null},P}\bigr)\,\lambda(dx).
\]
The clipping bound and cell masses imply
\[
\left|
\left(\int F_J(x)
\bigl(g_{T,P}(x)-c_{\mathrm{null},P}\bigr)\,\lambda(dx)\right)_j
\right|
\le\sqrt J\,\frac bJ=\frac b{\sqrt J}.
\]
Summing the squared coordinate bounds and using contraction by \(C_J\) gives
\[
\|\mu_P\|^2\le
\sum_{j=1}^J\frac{b^2}{J}=b^2.
\]
On \(\mathcal G_{n,P}\),
\[
W_P
\le b^2+16\sqrt\Lambda\,b+16\sqrt J\,\Lambda
\le2b^2+80\sqrt J\,\Lambda
\le2b^2+1024\sqrt J\,\Lambda.
\]
Here
\[
16\sqrt\Lambda\,b\le b^2+64\Lambda
\]
follows from \((b-8\sqrt\Lambda)^2\ge0\), and \(\sqrt J\ge1\).

The deterministic rule in \cref{obj:def:oracle-handle} rejects precisely when \(W_P\) exceeds its displayed cutoff. Its rejection event is therefore contained in \(\mathcal G_{n,P}^{\,c}\), and averaging the independent public randomization leaves that probability unchanged. Thus
\[
\mathsf R_n^{\mathrm e}(P,\phi^{*,\mathrm e}_{n,v})
\le0.01
\qquad(P\in H_0^{\mathrm{e,br}}(v)).
\]
% lean: CausalSmith.Stat.FinitepHomogeneityDensegamma.oracle_whole_null_size_le
% lean: CausalSmith.Stat.FinitepHomogeneityDensegamma.oracle_null_population_norm_le
% lean: CausalSmith.Stat.FinitepHomogeneityDensegamma.oracle_null_cutoff_le

\item For a general broader-model law \(P\), define its histogram coefficients and their associated constant by
\[
\theta_{\mathrm{hist},P}
=\int_{[0,1]}F_J(x)\tau_P(x)\,\lambda(dx),
\qquad
c_{\mathrm{hist},P}
=\frac1{\sqrt J}\sum_{j=1}^J(\theta_{\mathrm{hist},P})_j,
\qquad
\theta_{\mathrm{cen},P}=C_J\theta_{\mathrm{hist},P}.
\]
Subtracting the constant coefficient gives
\[
\theta_{\mathrm{cen},P}
=\int F_J(x)\bigl(\tau_P(x)-c_{\mathrm{hist},P}\bigr)\,\lambda(dx)
=\int F_J(x)\bigl((\Pi_J\tau_P)(x)-c_{\mathrm{hist},P}\bigr)\,\lambda(dx).
\]
The second equality holds because histogram projection preserves every cell integral.

Effect Hölder smoothness bounds the deviation from a cell average:
\[
|\tau_P(x)-(\Pi_J\tau_P)(x)|
\le20J^{-\gamma},
\qquad
\int(\tau_P-\Pi_J\tau_P)^2\,d\lambda
\le(20J^{-\gamma})^2.
\]
Indeed, every point in the same cell lies within distance \(1/J\), so averaging the Hölder bound over that cell proves the pointwise inequality.

The histogram projection is constant on each cell, and the residual has integral zero on each cell. Its residual is therefore orthogonal to the projection and to constants. Expanding the square, and using the histogram feature isometry, yields
\[
\int(\tau_P-c_{\mathrm{hist},P})^2\,d\lambda
=
\|\theta_{\mathrm{cen},P}\|^2
+\int(\tau_P-\Pi_J\tau_P)^2\,d\lambda.
\]
Also, expansion around the mean gives
\[
\int(\tau_P-c_{\mathrm{hist},P})^2\,d\lambda
=d(P)^2+(\bar\tau_P-c_{\mathrm{hist},P})^2
\ge d(P)^2.
\]
Consequently,
\[
d(P)^2
\le\|\theta_{\mathrm{cen},P}\|^2+(20J^{-\gamma})^2
\le\bigl(\|\theta_{\mathrm{cen},P}\|+20J^{-\gamma}\bigr)^2,
\]
and hence
\[
d(P)-20J^{-\gamma}\le\|\theta_{\mathrm{cen},P}\|.
\]

Finally,
\[
\mu_P-\theta_{\mathrm{cen},P}
=C_J\int F_J(x)\bigl(g_{T,P}(x)-\tau_P(x)\bigr)\,\lambda(dx).
\]
The same cellwise coefficient bound used for the null, now applied to
\(|g_{T,P}-\tau_P|\le b\), gives
\[
\|\mu_P-\theta_{\mathrm{cen},P}\|\le b.
\]
The reverse triangle inequality therefore proves the required population signal bound:
\[
d(P)-20J^{-\gamma}-b\le\|\mu_P\|.
\]
% lean: CausalSmith.Stat.FinitepHomogeneityDensegamma.centered_effect_histogram_distance
% lean: CausalSmith.Stat.FinitepHomogeneityDensegamma.oracle_effect_histogram_coefficient_distance
% lean: CausalSmith.Stat.FinitepHomogeneityDensegamma.oracle_population_clipping_error
% lean: CausalSmith.Stat.FinitepHomogeneityDensegamma.oracle_alternative_population_norm_ge

\item If
\[
d(P)\ge20J^{-\gamma}+5b+128\sqrt{\sqrt J\,\Lambda},
\]
the preceding population bound gives
\[
\|\mu_P\|\ge4b+128\sqrt{\sqrt J\,\Lambda}.
\]
Since \(J\ge1\),
\[
\sqrt\Lambda\le\sqrt{\sqrt J\,\Lambda},
\qquad
16\sqrt\Lambda\,\|\mu_P\|\le\frac{\|\mu_P\|^2}{2}.
\]
On \(\mathcal G_{n,P}\), it follows that
\[
\begin{aligned}
W_P
&\ge\|\mu_P\|^2
-16\sqrt\Lambda\,\|\mu_P\|-16\sqrt J\,\Lambda\\
&\ge\frac{\|\mu_P\|^2}{2}-16\sqrt J\,\Lambda\\
&\ge8b^2+8176\sqrt J\,\Lambda\\
&>2b^2+1024\sqrt J\,\Lambda.
\end{aligned}
\]
The penultimate inequality expands
\((4b+128\sqrt{\sqrt J\,\Lambda})^2/2\) and discards its nonnegative cross term. The last inequality is strict because \(\Lambda>0\). Thus the rule rejects throughout \(\mathcal G_{n,P}\), proving
\[
1-\mathsf R_n^{\mathrm e}(P,\phi^{*,\mathrm e}_{n,v})\le0.01
\]
at the displayed sufficient separation.
% lean: CausalSmith.Stat.FinitepHomogeneityDensegamma.oracle_power_of_mean_bound
% lean: CausalSmith.Stat.FinitepHomogeneityDensegamma.oracle_power_cutoff_lt

\item To express this separation in the public rate, define the variance-growth exponent
\[
t=\frac{2-p}{p-1}\ge0.
\]
Direct algebra gives
\[
2+t=\frac1q,
\qquad
E_0\left(2+t+\frac1{2\gamma}\right)=1.
\]
Upward dyadic rounding costs at most a factor of two:
\[
a^{-1/(p-1)}\le T\le2a^{-1/(p-1)},
\qquad
a^{-1/\gamma}\le J\le2a^{-1/\gamma}.
\]
The lower rounding bounds and the negative bias exponents imply
\[
b=20T^{1-p}\le20a,
\qquad
J^{-\gamma}\le a.
\]
Because \(0\le2-p\le1\), the upper rounding bounds imply
\[
T^{2-p}\le2a^{-t},
\qquad
\sqrt J\le\sqrt2\,a^{-1/(2\gamma)}.
\]
The exponent balance now gives
\[
s^{-1}a^{-t-1/(2\gamma)}=a^2,
\]
and consequently
\[
\sqrt J\,\Lambda
=\frac{80\sqrt J\,T^{2-p}}s
\le160\sqrt2\,a^2.
\]
Combining these estimates yields
\[
20J^{-\gamma}+5b+128\sqrt{\sqrt J\,\Lambda}
\le
\bigl(120+128\sqrt{160\sqrt2}\bigr)a.
\]

For every \(n\ge2\), the block size satisfies \(n\le3s\). Moreover,
\[
0<E_0=\frac{2\gamma q}{2\gamma+q}\le1,
\]
so
\[
a=s^{-E_0}
\le3^{E_0}n^{-E_0}
\le3\rho_n^{\mathrm e}(v).
\]
It follows that
\[
20J^{-\gamma}+5b+128\sqrt{\sqrt J\,\Lambda}
\le C_v^{\mathrm e}\rho_n^{\mathrm e}(v).
\]
Thus every broader-model law at or above the nominal upper separation has the asserted type-II bound, in particular under the stated gate
\(C_v^{\mathrm e}\rho_n^{\mathrm e}(v)<D_v^{\mathrm{e,br}}\).
All preceding estimates apply when \(s=1\) or \(J=1\); positivity of \(s,T,J,\Lambda\) suffices throughout.
% lean: CausalSmith.Stat.FinitepHomogeneityDensegamma.oracle_exponent_balance
% lean: CausalSmith.Stat.FinitepHomogeneityDensegamma.oracle_attaining_separation_le

\item Define
\[
r_{\mathrm{up}}=C_v^{\mathrm e}\rho_n^{\mathrm e}(v)>0.
\]
If \(r_{\mathrm{up}}<D_v^{\mathrm{e,br}}\), calibration makes
\(\phi^{*,\mathrm e}_{n,v}\) feasible at level \(1/10\), and its power bound gives
\[
B_n^{\mathrm{e,br}}(v,r_{\mathrm{up}})
\le
\sup_{\substack{P\in\mathcal M_v^{\mathrm{e,br}}\\d(P)\ge r_{\mathrm{up}}}}
\left\{1-\mathsf R_n^{\mathrm e}(P,\phi^{*,\mathrm e}_{n,v})\right\}
\le0.01\le\frac1{10}.
\]
For a nonempty alternative class, this supremum is bounded between zero and one; the empty-class convention assigns it zero and yields the same upper estimate. Here nonemptiness also follows from
\(r_{\mathrm{up}}<D_v^{\mathrm{e,br}}\): otherwise every model distance would be strictly below \(r_{\mathrm{up}}\), making \(r_{\mathrm{up}}\) an upper bound for their supremum. Thus \(r_{\mathrm{up}}\) belongs to the successful-radius set.

If \(D_v^{\mathrm{e,br}}\le r_{\mathrm{up}}\), the sentinel in
\cref{obj:def:oracle-broad-radius} instead gives
\[
r_n^{*,\mathrm{e,br}}(v)
\le D_v^{\mathrm{e,br}}\le r_{\mathrm{up}}.
\]
The defining set is nonempty because it contains the sentinel, and it is bounded below by zero. In both cases,
\[
r_n^{*,\mathrm{e,br}}(v)\le r_{\mathrm{up}}.
\]
Together with the lower bound, this proves the radius inequalities. The witnesses constructed above have all the required model memberships, common function input, distance bounds, normalization, total variation comparison, and positive-weight error witness.
% lean: CausalSmith.Stat.FinitepHomogeneityDensegamma.oracleBroadCriticalRadius_le_of_test
% lean: CausalSmith.Stat.FinitepHomogeneityDensegamma.oracle_frontier_broad_class

\item For \(n\ge N_v^{\mathrm e}\), the ceiling definition implies
\[
n\ge\left(\frac{2C_v^{\mathrm e}}{d_0}\right)^{1/E_0}.
\]
The base and \(E_0\) are positive. Raising this inequality to the negative power \(-E_0\) therefore gives
\[
C_v^{\mathrm e}n^{-E_0}
\le
C_v^{\mathrm e}
\left[
\left(\frac{2C_v^{\mathrm e}}{d_0}\right)^{1/E_0}
\right]^{-E_0}
=
C_v^{\mathrm e}\left(\frac{2C_v^{\mathrm e}}{d_0}\right)^{-1}
=\frac{d_0}{2}.
\]
% lean: CausalSmith.Stat.FinitepHomogeneityDensegamma.oracle_attainment_threshold

\item For \(P\in\mathcal M_v^{\mathrm{e,br}}\), the conditional arm-mean identities hold simultaneously for \(\lambda\)-almost every \(x\):
\[
\int_{\mathbb R}y\,Q_{0,P}(dy\mid x)=m_{0,P}(x),
\qquad
\int_{\mathbb R}y\,Q_{1,P}(dy\mid x)
=m_{0,P}(x)+\tau_P(x).
\]
Overlap gives \(e_P(x)>0\) and \(1-e_P(x)>0\). Pulling the constant denominators out of the integrals and cancelling them yields
\[
\begin{aligned}
&e_P(x)\int_{\mathbb R}\frac{y}{e_P(x)}\,Q_{1,P}(dy\mid x)
+(1-e_P(x))\int_{\mathbb R}\frac{-y}{1-e_P(x)}\,Q_{0,P}(dy\mid x)\\
&\qquad=
\int_{\mathbb R}y\,Q_{1,P}(dy\mid x)
-\int_{\mathbb R}y\,Q_{0,P}(dy\mid x)
=\tau_P(x).
\end{aligned}
\]
The conditional first moments are finite almost everywhere by the raw moment bound.
% lean: CausalSmith.Stat.FinitepHomogeneityDensegamma.oracle_untruncated_mean_identity

\item Finally, the coordinate maps \(v\mapsto p\) and \(v\mapsto\gamma\) are continuous in the product topology. On the admissible domain,
\[
p>0,\qquad
q=\frac{p-1}{p}>0,\qquad
2\gamma+q>0.
\]
Hence \(q\) and \(E_0\) are continuous there, and \(E_0\) is everywhere positive.

For a compact subset \(\mathcal K\) of that domain, define
\[
E_{\min,\mathcal K}
=
\begin{cases}
\displaystyle\min_{v\in\mathcal K}E_0(v),&\mathcal K\ne\varnothing,\\[1mm]
1,&\mathcal K=\varnothing.
\end{cases}
\]
In the nonempty case, continuity and compactness supply an attained minimum, whose value is positive by admissibility. In the empty case the assertions are vacuous. Thus
\[
E_{\min,\mathcal K}>0,
\qquad
E_{\min,\mathcal K}\le E_0(v)\quad(v\in\mathcal K).
\]
Together with the global lower-multiplier bound and the fixed finite upper multiplier,
\[
c_v^{\mathrm e}\ge\frac1{64\sqrt3},
\qquad
C_v^{\mathrm e}=3\bigl(120+128\sqrt{160\sqrt2}\bigr),
\]
this proves the asserted local uniformity.
% lean: CausalSmith.Stat.FinitepHomogeneityDensegamma.oracle_compact_uniformity
\end{enumerate}
\end{proof}

\begin{proof}[Proof of \cref{obj:thm:oracle-frontier-resolved}]
Fix \(v\in\mathcal V\). Admissibility gives
\[
1<p\le2,\qquad \frac14\le\gamma\le1,\qquad
0<q=\frac{p-1}{p}\le\frac12,\qquad
2\gamma+q>0,\qquad E_0>0.
\]
We use the constants and scales in the statement.

\begin{enumerate}
\item Since \(-1\le-\gamma\) and the base \(4\) exceeds one,
\[
4^{-1}\le4^{-\gamma},
\qquad
\frac{1}{64\sqrt3}
=\frac{4^{-1}}{16\sqrt3}
\le c^{\mathrm e}_v.
\]
Also,
\[
C^{\mathrm e}_v
=3\left(120+128\sqrt{160\sqrt2}\right)>0.
\]
% lean: CausalSmith.Stat.FinitepHomogeneityDensegamma.cOracle_uniform_lower
% lean: CausalSmith.Stat.FinitepHomogeneityDensegamma.COr_pos

\item Fix an integer \(n\ge2\), and define the lower separation by
\[
\varrho_-=c^{\mathrm e}_v\rho_n^{\mathrm e}(v).
\]
By \cref{obj:lem:paired-tent-full-record-lower}, there are a null law \(P_0\in H_0(v)\) and a finite prior
\[
\pi_1=\sum_{j=1}^{k}\omega_j\delta_{P_j},
\qquad
\omega_j\ge0,\qquad
\sum_{j=1}^{k}\omega_j=1,
\qquad
\exists j:\omega_j>0.
\]
For every positive-weight support law, that result supplies
\[
P_j\in\mathcal M_v,\qquad
e_{P_j}(x)=\frac12=e_{P_0}(x)\quad(x\in[0,1]),
\qquad
\varrho_-\le d(P_j)<d_0\le D_v.
\]
It also supplies
\[
\mathbb P_{\pi_1,n}
=\sum_{j=1}^{k}\omega_jP_j^{\otimes n},
\qquad
\operatorname{TV}\!\left(P_0^{\otimes n},\mathbb P_{\pi_1,n}\right)<\frac12,
\qquad
B_n^{\mathrm e}(v,\varrho_-)\ge\frac25.
\]
% lean: CausalSmith.Stat.FinitepHomogeneityDensegamma.paired_tent_full_record_lower

Choose an index \(j_+\in\{1,\ldots,k\}\) with \(\omega_{j_+}>0\). Positivity of the lower multiplier and of \(n^{-E_0}\) yields
\[
0<\varrho_-\le d(P_{j_+})<D_v.
\]
Thus the alternatives at separation \(\varrho_-\) are nonempty. For every \(0<\varrho\le\varrho_-\), their inclusion in the alternatives at separation \(\varrho\) gives
\[
B_n^{\mathrm e}(v,\varrho_-)
\le B_n^{\mathrm e}(v,\varrho).
\]
The size constraint is the same in both infima. Consequently, every successful separation below \(D_v\) in \cref{obj:def:oracle-radius} is at least \(\varrho_-\), because
\[
B_n^{\mathrm e}(v,\varrho)\ge\frac25>\frac1{10}
\qquad(0<\varrho\le\varrho_-).
\]
The additional candidate \(D_v\) also exceeds \(\varrho_-\). Taking the infimum therefore gives
\[
c^{\mathrm e}_v\rho_n^{\mathrm e}(v)
=\varrho_-\le r_n^{*,\mathrm e}(v).
\]
% lean: CausalSmith.Stat.FinitepHomogeneityDensegamma.oracleCriticalRadius_ge_of_failed

\item We next establish the null calibration of the rule in \cref{obj:def:oracle-handle}. Use its tuning quantities
\[
s=\lfloor n/2\rfloor,\qquad a=s^{-E_0},\qquad
T=2^{\lceil\log_2(a^{-1/(p-1)})\rceil},\qquad
J=2^{\lceil\log_2(a^{-1/\gamma})\rceil},
\]
\[
b=20T^{1-p},\qquad
\Lambda=\frac{80T^{2-p}}s.
\]
Here \(s\ge1\), \(0<a\le1\), \(T\ge1\), \(J\ge1\), \(b\ge0\), and \(\Lambda>0\).
With \(F_J\) from \cref{obj:def:projection}, define
\[
u=J^{-1/2}(1,\ldots,1)^T,\qquad
C_J=I_J-uu^T.
\]
The cells have Lebesgue measure \(1/J\), so
\[
\int_0^1F_J(x)F_J(x)^T\,d\lambda(x)=I_J,
\qquad
\int_0^1F_J(x)\,d\lambda(x)=u,
\qquad
C_JF_J(x)=F_J(x)-u.
\]
Thus \(C_J\) is a symmetric orthogonal projection and
\[
\|C_J\xi_{\mathrm{dir}}\|_2\le\|\xi_{\mathrm{dir}}\|_2
\qquad(\xi_{\mathrm{dir}}\in\mathbb R^J).
\]

For \(P\in\mathcal M_v\), define the clipped conditional mean and second moment on \([0,1]\) by
\[
g_{T,P}(x)
=\int_{\mathbb R}\ell_T(y)\,Q_{1,P}(dy\mid x)
-\int_{\mathbb R}\ell_T(y)\,Q_{0,P}(dy\mid x),
\]
\[
m_{2,T,P}(x)
=\frac{\int_{\mathbb R}\ell_T(y)^2\,Q_{1,P}(dy\mid x)}{e_P(x)}
+\frac{\int_{\mathbb R}\ell_T(y)^2\,Q_{0,P}(dy\mid x)}{1-e_P(x)}.
\]
These functions are finite and measurable, since clipping is bounded and overlap makes both denominators positive.

For every \(y\in\mathbb R\),
\[
|y-\ell_T(y)|\le |y|^pT^{1-p},
\qquad
\ell_T(y)^2\le |y|^pT^{2-p}.
\]
Indeed, on \(|y|\le T\), the clipping remainder vanishes and
\(|y|^{2-p}\le T^{2-p}\); on \(|y|>T\), the remainder is at most \(|y|\) and the clipped square equals \(T^2\). The moment restriction in \cref{obj:ass:raw-moment} therefore gives, for each \(a'\in\{0,1\}\) and almost every \(x\),
\[
\left|\int_{\mathbb R}(y-\ell_T(y))\,Q_{a',P}(dy\mid x)\right|
\le10T^{1-p},
\qquad
\int_{\mathbb R}\ell_T(y)^2\,Q_{a',P}(dy\mid x)
\le10T^{2-p}.
\]
The same moment restriction ensures integrability of \(y\), using
\(|y|\le1+|y|^p\).
The arm-mean identities now yield
\[
g_{T,P}(x)-\tau_P(x)
=
\int_{\mathbb R}(y-\ell_T(y))\,Q_{0,P}(dy\mid x)
-\int_{\mathbb R}(y-\ell_T(y))\,Q_{1,P}(dy\mid x),
\]
and hence
\[
|g_{T,P}(x)-\tau_P(x)|\le b,
\qquad
m_{2,T,P}(x)
\le10T^{2-p}
\left(\frac1{e_P(x)}+\frac1{1-e_P(x)}\right)
\le80T^{2-p}
\]
almost everywhere. These are respectively the conditional mean and second moment bounds for \(R_i\), since the supplied true propensity agrees with its overlap-clipped version.
% lean: CausalSmith.Stat.FinitepHomogeneityDensegamma.oracle_clipping_moments

Define the common block mean and the two histogram coefficient vectors by
\[
\mu_P=\mathbb E_{P^{\otimes n}}[V_1]
=\mathbb E_{P^{\otimes n}}[V_2]
=C_J\theta_{T,P},
\]
\[
\theta_{T,P}=\int_0^1g_{T,P}(x)F_J(x)\,d\lambda(x),
\qquad
\theta_{\tau,P}=\int_0^1\tau_P(x)F_J(x)\,d\lambda(x).
\]
The equality of the block means follows by averaging \(s\) identically distributed records in each block.

For a null law, define its constant effect by
\[
c_P=\tau_P(0),\qquad \tau_P(x)=c_P\quad(x\in[0,1]).
\]
Because \(C_Ju=0\),
\[
\mu_P=C_J\int_0^1(g_{T,P}(x)-c_P)F_J(x)\,d\lambda(x).
\]
For every histogram coordinate \(j\),
\[
\left|
\left(\int_0^1(g_{T,P}-c_P)F_J\,d\lambda\right)_j
\right|
=
\left|\sqrt J\int_{I_{j,J}}(g_{T,P}-c_P)\,d\lambda\right|
\le\frac b{\sqrt J}.
\]
Summing the squared coordinate bounds and applying the contraction \(C_J\) gives
\[
\|\mu_P\|_2^2
\le\sum_{j=1}^{J}\frac{b^2}{J}=b^2,
\qquad
\|\mu_P\|_2\le b.
\]
% lean: CausalSmith.Stat.FinitepHomogeneityDensegamma.oracle_null_population_norm_le

For any model law and any \(\xi_{\mathrm{dir}}\in\mathbb R^J\), the conditional second moment bound and the feature isometry give
\[
\mathbb E_P\!\left[
\left\langle\xi_{\mathrm{dir}},C_JF_J(X)R\right\rangle^2
\right]
\le80T^{2-p}\int_0^1
\left\langle C_J\xi_{\mathrm{dir}},F_J(x)\right\rangle^2\,d\lambda(x)
\le80T^{2-p}\|\xi_{\mathrm{dir}}\|_2^2.
\]
Here the single-record score is defined by
\[
R=\frac{A\ell_T(Y)}{e_P(X)}
-\frac{(1-A)\ell_T(Y)}{1-e_P(X)}.
\]
Independence within a block gives, for \(b'\in\{1,2\}\),
\[
\operatorname{Var}_{P^{\otimes n}}
\!\left(\left\langle\xi_{\mathrm{dir}},V_{b'}\right\rangle\right)
=
\frac1s\operatorname{Var}_P
\!\left(\left\langle\xi_{\mathrm{dir}},C_JF_J(X)R\right\rangle\right)
\le\Lambda\|\xi_{\mathrm{dir}}\|_2^2.
\]
% lean: CausalSmith.Stat.FinitepHomogeneityDensegamma.oracleBlockVec_variance_le

Define
\[
\Delta_{b',P}=V_{b'}-\mu_P\quad(b'\in\{1,2\}),
\qquad
W_{\mathrm e}=\langle V_1,V_2\rangle.
\]
The centered vectors are independent and have mean zero. In particular,
\[
\mathbb E_{P^{\otimes n}}[W_{\mathrm e}]=\|\mu_P\|_2^2,
\]
\[
W_{\mathrm e}-\|\mu_P\|_2^2
=
\langle\mu_P,\Delta_{1,P}\rangle
+\langle\mu_P,\Delta_{2,P}\rangle
+\langle\Delta_{1,P},\Delta_{2,P}\rangle.
\]
The expectations of all three cross products in the square of this sum vanish: the product of the two linear terms factors into their zero means, and conditioning on the other block makes each linear--bilinear cross product zero. To bound the remaining bilinear square, let
\[
\mathbf e^{(j)}
=\bigl(\mathbf1_{\{j'=j\}}\bigr)_{j'=1}^{J}
\qquad(j=1,\ldots,J)
\]
be the coordinate unit vectors. Independence and the directional second moment bounds imply
\[
\mathbb E_{P^{\otimes n}}
\!\left[\langle\Delta_{1,P},\Delta_{2,P}\rangle^2\right]
\le\Lambda\,\mathbb E_{P^{\otimes n}}\|\Delta_{2,P}\|_2^2
=\Lambda\sum_{j=1}^{J}
\mathbb E_{P^{\otimes n}}
\!\left[\langle\mathbf e^{(j)},\Delta_{2,P}\rangle^2\right]
\le J\Lambda^2.
\]
Consequently,
\[
\operatorname{Var}_{P^{\otimes n}}(W_{\mathrm e})
\le2\Lambda\|\mu_P\|_2^2+J\Lambda^2.
\]
All the required moments are finite because the clipped score and the finite-rank features are bounded.
% lean: CausalSmith.Stat.FinitepHomogeneityDensegamma.oracleBlockVec_quadratic_variance_le

The squared-deviation inequality now gives
\[
\begin{aligned}
&P^{\otimes n}\!\left(
|W_{\mathrm e}-\|\mu_P\|_2^2|
>16\{\sqrt\Lambda\|\mu_P\|_2+\sqrt J\,\Lambda\}
\right)\\
&\quad\le
\frac{2\Lambda\|\mu_P\|_2^2+J\Lambda^2}
{256\{\sqrt\Lambda\|\mu_P\|_2+\sqrt J\,\Lambda\}^2}
\le\frac{2}{256}<0.01.
\end{aligned}
\]
The denominator is positive since \(J\ge1\) and \(\Lambda>0\).
Define the good event by
\[
\mathcal G_P=
\left\{
|W_{\mathrm e}-\|\mu_P\|_2^2|
\le16\{\sqrt\Lambda\|\mu_P\|_2+\sqrt J\,\Lambda\}
\right\}.
\]
Thus \(P^{\otimes n}(\mathcal G_P^c)\le0.01\).
% lean: CausalSmith.Stat.FinitepHomogeneityDensegamma.oracleBlockVec_tail_le

For a null law, on \(\mathcal G_P\) we have
\[
W_{\mathrm e}
\le b^2+16\sqrt\Lambda\,b+16\sqrt J\,\Lambda
\le2b^2+80\sqrt J\,\Lambda
\le2b^2+1024\sqrt J\,\Lambda.
\]
The middle inequality follows from
\[
16\sqrt\Lambda\,b\le b^2+64\Lambda,
\qquad
\Lambda\le\sqrt J\,\Lambda,
\]
the first being equivalent to \((b-8\sqrt\Lambda)^2\ge0\).
The rejection event is therefore contained in \(\mathcal G_P^c\). Since the rule ignores the public randomization variable,
\[
\mathsf R_n^{\mathrm e}(P,\phi^{*,\mathrm e}_{n,v})
=
P^{\otimes n}\!\left(
W_{\mathrm e}>2b^2+1024\sqrt J\,\Lambda
\right)
\le0.01
\qquad(P\in H_0(v)).
\]
These calculations also apply when \(s=1\); when \(J=1\), \(C_J=0\) and both centered block vectors vanish.
% lean: CausalSmith.Stat.FinitepHomogeneityDensegamma.oracle_whole_null_size_le

\item We establish the alternative population bound before applying the same concentration event. Define
\[
\overline\tau_P=\int_0^1\tau_P(x)\,d\lambda(x),
\qquad
c_{P,J}=\frac1{\sqrt J}\sum_{j=1}^{J}(\theta_{\tau,P})_j,
\qquad
h_{P,J}(x)=\tau_P(x)-c_{P,J}.
\]
The coefficient formula for constants and the definition of \(C_J\) give
\[
C_J\theta_{\tau,P}
=\int_0^1h_{P,J}(x)F_J(x)\,d\lambda(x).
\]
On each cell, the histogram projection is the cell average. Hence
\[
\tau_P(x)-(\Pi_J\tau_P)(x)
=J\int_{I_{j,J}}\{\tau_P(x)-\tau_P(z)\}\,d\lambda(z)
\qquad(x\in I_{j,J}),
\]
and \cref{obj:ass:effect-smoothness} gives
\[
|\tau_P(x)-(\Pi_J\tau_P)(x)|\le20J^{-\gamma},
\qquad
\int_0^1\{\tau_P-\Pi_J\tau_P\}^2\,d\lambda
\le(20J^{-\gamma})^2.
\]
Subtracting a constant commutes with \(\Pi_J\). Its cellwise residual has integral zero on every cell, and is therefore orthogonal to the cellwise constant projection. The feature isometry then yields
\[
\int_0^1h_{P,J}^2\,d\lambda
=
\|C_J\theta_{\tau,P}\|_2^2
+\int_0^1\{\tau_P-\Pi_J\tau_P\}^2\,d\lambda.
\]
Expanding the distance to the constant \(c_{P,J}\) also gives
\[
\int_0^1h_{P,J}^2\,d\lambda
=d(P)^2+(\overline\tau_P-c_{P,J})^2
\ge d(P)^2.
\]
It follows that
\[
d(P)^2
\le\|C_J\theta_{\tau,P}\|_2^2+(20J^{-\gamma})^2
\le\bigl(\|C_J\theta_{\tau,P}\|_2+20J^{-\gamma}\bigr)^2,
\]
so
\[
d(P)-20J^{-\gamma}\le\|C_J\theta_{\tau,P}\|_2.
\]
% lean: CausalSmith.Stat.FinitepHomogeneityDensegamma.centered_effect_histogram_distance

The same cellwise coefficient bound used for the null, now applied to \(g_{T,P}-\tau_P\), gives
\[
\|\mu_P-C_J\theta_{\tau,P}\|_2
=
\left\|C_J\int_0^1(g_{T,P}-\tau_P)F_J\,d\lambda\right\|_2
\le b.
\]
Combining this with the preceding lower bound and the reverse triangle inequality yields
\[
d(P)-20J^{-\gamma}-b\le\|\mu_P\|_2.
\]
% lean: CausalSmith.Stat.FinitepHomogeneityDensegamma.oracle_alternative_population_norm_ge

Define the noise scale by
\[
\zeta_{\mathrm{noise}}=\sqrt{\sqrt J\,\Lambda}>0.
\]
If
\[
d(P)\ge20J^{-\gamma}+5b+128\zeta_{\mathrm{noise}},
\]
then
\[
\|\mu_P\|_2\ge4b+128\zeta_{\mathrm{noise}}.
\]
Since \(\sqrt\Lambda\le\zeta_{\mathrm{noise}}\), this implies
\[
16\sqrt\Lambda\,\|\mu_P\|_2\le\frac{\|\mu_P\|_2^2}{2}.
\]
On \(\mathcal G_P\), therefore,
\[
\begin{aligned}
W_{\mathrm e}
&\ge\|\mu_P\|_2^2
-16\sqrt\Lambda\,\|\mu_P\|_2
-16\zeta_{\mathrm{noise}}^2\\
&\ge\frac{\|\mu_P\|_2^2}{2}-16\zeta_{\mathrm{noise}}^2\\
&\ge8b^2+8176\zeta_{\mathrm{noise}}^2
>2b^2+1024\zeta_{\mathrm{noise}}^2.
\end{aligned}
\]
Here the third inequality uses
\[
\frac{(4b+128\zeta_{\mathrm{noise}})^2}{2}
\ge8b^2+8192\zeta_{\mathrm{noise}}^2,
\]
and the last is strict because \(\zeta_{\mathrm{noise}}>0\).
Thus failure to reject is contained in \(\mathcal G_P^c\), proving
\[
1-\mathsf R_n^{\mathrm e}(P,\phi^{*,\mathrm e}_{n,v})\le0.01
\]
under the displayed separation condition.
% lean: CausalSmith.Stat.FinitepHomogeneityDensegamma.oracle_power_of_mean_bound

\item We now bound that sufficient separation by the scale in the statement. Define the variance-growth exponent by
\[
t=\frac{2-p}{p-1}\ge0.
\]
Both dyadic targets are at least one, and upward rounding gives
\[
a^{-1/(p-1)}\le T\le2a^{-1/(p-1)},
\qquad
a^{-1/\gamma}\le J\le2a^{-1/\gamma}.
\]
Since \(1-p<0\) and \(-\gamma<0\),
\[
b=20T^{1-p}\le20a,
\qquad
J^{-\gamma}\le a.
\]
Also \(0\le2-p\le1\), so
\[
T^{2-p}
\le2^{2-p}a^{-t}\le2a^{-t},
\qquad
\sqrt J\le\sqrt2\,a^{-1/(2\gamma)}.
\]
Substitution of the definitions of \(q\) and \(E_0\) gives the balance identity
\[
E_0\left(2+t+\frac1{2\gamma}\right)=1.
\]
Because \(a=s^{-E_0}\), this identity implies
\[
\frac{a^{-t-1/(2\gamma)}}s=a^2.
\]
Consequently,
\[
\sqrt J\,\Lambda
\le160\sqrt2\,\frac{a^{-t-1/(2\gamma)}}s
=160\sqrt2\,a^2,
\qquad
\zeta_{\mathrm{noise}}\le\sqrt{160\sqrt2}\,a.
\]

For every \(n\ge2\), \(n\le3s\). Moreover,
\[
2\gamma q\le2\gamma+q,\qquad 0<E_0\le1,
\qquad 3^{E_0}\le3.
\]
Monotonicity of the negative power therefore gives
\[
a=s^{-E_0}
\le3^{E_0}n^{-E_0}
\le3\rho_n^{\mathrm e}(v).
\]
Combining the approximation, clipping, and covariance bounds yields
\[
\begin{aligned}
20J^{-\gamma}+5b+128\zeta_{\mathrm{noise}}
&\le\left(120+128\sqrt{160\sqrt2}\right)a\\
&\le C^{\mathrm e}_v\rho_n^{\mathrm e}(v).
\end{aligned}
\]
Thus, whenever \(C^{\mathrm e}_v\rho_n^{\mathrm e}(v)<D_v\), every model law at distance at least \(C^{\mathrm e}_v\rho_n^{\mathrm e}(v)\) has type-II error at most \(0.01\).
% lean: CausalSmith.Stat.FinitepHomogeneityDensegamma.oracle_attaining_separation_le

\item Define the upper separation by
\[
\varrho_+=C^{\mathrm e}_v\rho_n^{\mathrm e}(v)>0.
\]
The witnesses already give \(D_v>0\). If \(\varrho_+<D_v\), the calibrated rule satisfies the size constraint \(1/10\), and its alternative errors give
\[
\sup_{P\in\mathcal M_v:\,d(P)\ge\varrho_+}
\left\{1-\mathsf R_n^{\mathrm e}(P,\phi^{*,\mathrm e}_{n,v})\right\}
\le0.01\le\frac1{10}.
\]
For an empty alternative index set, the real-valued supremum is assigned the value \(0\), so the same bound applies in that case. All worst-error suprema are nonnegative, and the admissible test family contains the calibrated rule. Taking its value in the infimum of \cref{obj:def:oracle-risk} therefore shows
\[
B_n^{\mathrm e}(v,\varrho_+)\le\frac1{10}.
\]
Hence \(\varrho_+\) is a candidate in \cref{obj:def:oracle-radius}, and
\[
r_n^{*,\mathrm e}(v)\le\varrho_+.
\]
If instead \(D_v\le\varrho_+\), the candidate \(D_v\) gives
\[
r_n^{*,\mathrm e}(v)\le D_v\le\varrho_+.
\]
Both cases establish the upper radius bound.
% lean: CausalSmith.Stat.FinitepHomogeneityDensegamma.oracleCriticalRadius_le_of_test

\item Retain the prior and null law already constructed, and let \(\phi^{\mathrm e}\) satisfy the stipulated size bound on \(H_0(v)\). Define its seed-averaged map with the common propensity supplied by
\[
\overline\phi^{\mathrm e}_{P_0}(\mathcal D_n)
=\int_0^1\phi^{\mathrm e}(e_{P_0},\mathcal D_n,u_{\mathrm{seed}})\,d\lambda(u_{\mathrm{seed}}),
\qquad
0\le\overline\phi^{\mathrm e}_{P_0}\le1.
\]
This is a measurable function on the complete-record sample space. Integrating event-probability differences over its level sets gives
\[
\left|
\int\overline\phi^{\mathrm e}_{P_0}\,d\mathbb P_{\pi_1,n}
-\int\overline\phi^{\mathrm e}_{P_0}\,dP_0^{\otimes n}
\right|
\le\operatorname{TV}\!\left(P_0^{\otimes n},\mathbb P_{\pi_1,n}\right).
\]
Indeed, each integral equals the integral over \(\vartheta_{\mathrm{lev}}\in[0,1]\) of the probability of the measurable level set
\(\{\overline\phi^{\mathrm e}_{P_0}>\vartheta_{\mathrm{lev}}\}\).

Every positive-weight support law has propensity \(e_{P_0}\). Finite linearity of integration and the null size bound thus yield
\[
\sum_{j=1}^{k}\omega_j
\mathsf R_n^{\mathrm e}(P_j,\phi^{\mathrm e})
\le
\mathsf R_n^{\mathrm e}(P_0,\phi^{\mathrm e})
+\operatorname{TV}\!\left(P_0^{\otimes n},\mathbb P_{\pi_1,n}\right)
<\frac1{10}+\frac12=\frac35.
\]
If every positive-weight index had type-II error below \(2/5\), its rejection probability would exceed \(3/5\). The zero-weight terms vanish, and normalization would then imply
\[
\sum_{j=1}^{k}\omega_j
\mathsf R_n^{\mathrm e}(P_j,\phi^{\mathrm e})
\ge\frac35\sum_{j=1}^{k}\omega_j=\frac35,
\]
a contradiction. Therefore
\[
\exists j:\quad
\omega_j>0,\qquad
1-\mathsf R_n^{\mathrm e}(P_j,\phi^{\mathrm e})\ge\frac25.
\]
This proves the asserted individual error witness for the same common-propensity prior.
% lean: CausalSmith.Stat.FinitepHomogeneityDensegamma.oracle_single_null_prior_error_witness

\item For \(n\ge N^{\mathrm e}_v\), the definition of the ceiling gives
\[
n\ge
\max\left\{2,\left(\frac{2C^{\mathrm e}_v}{d_0}\right)^{1/E_0}\right\}.
\]
Since \(E_0>0\), raising the lower bound for \(n\) to the power \(-E_0\) reverses the inequality:
\[
n^{-E_0}
\le
\left\{
\left(\frac{2C^{\mathrm e}_v}{d_0}\right)^{1/E_0}
\right\}^{-E_0}
=\left(\frac{2C^{\mathrm e}_v}{d_0}\right)^{-1}.
\]
Multiplication by \(C^{\mathrm e}_v>0\) proves
\[
C^{\mathrm e}_v\rho_n^{\mathrm e}(v)\le\frac{d_0}{2}.
\]
% lean: CausalSmith.Stat.FinitepHomogeneityDensegamma.oracle_attainment_threshold

\item Let \(P\in\mathcal M_v\). On the common almost-everywhere set where the two conditional arm-mean identities hold,
\[
\int_{\mathbb R}y\,Q_{0,P}(dy\mid x)=m_{0,P}(x),
\qquad
\int_{\mathbb R}y\,Q_{1,P}(dy\mid x)=m_{0,P}(x)+\tau_P(x).
\]
Overlap gives \(e_P(x)>0\) and \(1-e_P(x)>0\) at every \(x\). Pulling the constant denominators outside the arm integrals and cancelling them therefore yields
\[
\begin{aligned}
&e_P(x)\int_{\mathbb R}\frac{y}{e_P(x)}\,Q_{1,P}(dy\mid x)
+(1-e_P(x))
\int_{\mathbb R}\frac{-y}{1-e_P(x)}\,Q_{0,P}(dy\mid x)\\
&\qquad=
\int_{\mathbb R}y\,Q_{1,P}(dy\mid x)
-\int_{\mathbb R}y\,Q_{0,P}(dy\mid x)
=\tau_P(x).
\end{aligned}
\]
This is the claimed identity for the untruncated score.
% lean: CausalSmith.Stat.FinitepHomogeneityDensegamma.oracle_untruncated_mean_identity

\item Finally, let \(\mathcal K\subseteq\mathcal V\) be compact. The coordinate functions \(p\) and \(\gamma\) are continuous in the inherited topology. Since \(p>1\) and \(2\gamma+(p-1)/p>0\), the function
\[
v\longmapsto
E_0(v)=
\frac{2\gamma\,((p-1)/p)}
{2\gamma+(p-1)/p}
\]
is continuous and positive throughout \(\mathcal V\). If \(\mathcal K\) is nonempty, it attains a positive minimum; if \(\mathcal K\) is empty, choose the value one. Thus there is
\[
\eta_{\mathcal K}>0,
\qquad
\eta_{\mathcal K}\le E_0(v)\quad(v\in\mathcal K).
\]
The multiplier bounds already proved give
\[
c^{\mathrm e}_v\ge\frac1{64\sqrt3},
\qquad
C^{\mathrm e}_v=3\left(120+128\sqrt{160\sqrt2}\right)
\qquad(v\in\mathcal K).
\]
Using \(d_0=1/(8\sqrt2)\) from \cref{obj:lem:paired-tent-full-record-lower}, the base \(2C^{\mathrm e}_v/d_0\) exceeds one and is independent of \(v\). Therefore
\[
N^{\mathrm e}_v
\le
\left\lceil
\max\left\{
2,\left(\frac{2C^{\mathrm e}_v}{d_0}\right)^{1/\eta_{\mathcal K}}
\right\}
\right\rceil
\qquad(v\in\mathcal K).
\]
This supplies a common finite threshold bound. Positivity of the denominators and continuity of the exponent also hold at \(p=2\) and \(\gamma=1/4\), so the same compact-uniform argument includes those faces.
% lean: CausalSmith.Stat.FinitepHomogeneityDensegamma.oracle_compact_uniformity
\end{enumerate}
\end{proof}

\begin{proof}[Proof of \cref{obj:thm:exact-propensity-preservation-boundary}]
\begin{enumerate}
\item Fix \(v=(p,\alpha,\beta,\gamma)\in\mathcal V\), and use the quantities of \cref{obj:def:sharp-frontier-scales}. By \cref{obj:thm:heavy-tail-comparison,obj:thm:oracle-frontier-resolved},
\[
c_v>0,\qquad C_v>0,\qquad
c^{\mathrm e}_v\ge\frac{1}{64\sqrt3}>0,\qquad
C^{\mathrm e}_v>0.
\]
Consequently,
\[
\frac{c_v}{C^{\mathrm e}_v}>0,
\qquad
\frac{C_v}{c^{\mathrm e}_v}>0.
\]
% lean: CausalSmith.Stat.FinitepHomogeneityDensegamma.exact_propensity_preservation_boundary

\item We establish the identities for the comparison exponent. Admissibility gives
\[
p-1>0,\qquad p>0,\qquad q>0,\qquad
\gamma>0,\qquad D_p>0,
\]
and hence
\[
\frac{2\gamma q}{(2\gamma+q)D_p}>0.
\]
The phase identity and channel selection in \cref{obj:thm:heavy-tail-comparison} give
\[
E_4(p)-E_0(p)
=
\frac{2\gamma q}{(2\gamma+q)D_p}\bigl(F(p)-1\bigr),
\]
\[
F(p)\ge1\ \Longrightarrow\ E(p)=E_0(p),
\qquad
F(p)<1\ \Longrightarrow\ E(p)=E_4(p).
\]

If \(F(p)\ge1\), the positive prefactor implies
\[
E_0(p)-E_4(p)\le0.
\]
Thus, using the definition of \(G(v)\),
\[
G(v)=0=E_0(p)-E(p)
=
\frac{2\gamma q}{(2\gamma+q)D_p}\max\{0,1-F(p)\}.
\]
If \(F(p)<1\), the same identity implies
\[
E_0(p)-E_4(p)
=
\frac{2\gamma q}{(2\gamma+q)D_p}\bigl(1-F(p)\bigr)>0.
\]
Channel selection now gives
\[
G(v)=E_0(p)-E_4(p)=E_0(p)-E(p)
=
\frac{2\gamma q}{(2\gamma+q)D_p}\max\{0,1-F(p)\}.
\]
Together, these cases establish
\[
G(v)=E_0(p)-E(p),
\qquad
G(v)=
\frac{2\gamma q}{(2\gamma+q)D_p}\max\{0,1-F(p)\}.
\]
In particular,
\[
F(p)\ge1\ \Longrightarrow\ G(v)=0.
\]
% lean: CausalSmith.Stat.FinitepHomogeneityDensegamma.preservation_gap_identity

\item For every integer \(n\ge2\), the two frontier conclusions in
\cref{obj:thm:heavy-tail-comparison,obj:thm:oracle-frontier-resolved} supply
\[
c_v\rho_n(v)\le r_n^*(v)\le C_v\rho_n(v),
\qquad
c^{\mathrm e}_v\rho_n^{\mathrm e}(v)
\le r_n^{*,\mathrm e}(v)
\le C^{\mathrm e}_v\rho_n^{\mathrm e}(v).
\]
Since \(n>0\), both power scales are positive:
\[
\rho_n(v)=n^{-E(p)}>0,
\qquad
\rho_n^{\mathrm e}(v)=n^{-E_0(p)}>0.
\]
The lower bounds therefore imply
\[
r_n^*(v)>0,\qquad r_n^{*,\mathrm e}(v)>0.
\]
Dividing the original-radius lower bound by the supplied-radius upper bound, and the original-radius upper bound by the supplied-radius lower bound, yields
\[
\frac{c_v\rho_n(v)}
     {C^{\mathrm e}_v\rho_n^{\mathrm e}(v)}
\le
\frac{r_n^*(v)}{r_n^{*,\mathrm e}(v)}
\le
\frac{C_v\rho_n(v)}
     {c^{\mathrm e}_v\rho_n^{\mathrm e}(v)}.
\]
The scale quotient satisfies
\[
\frac{\rho_n(v)}{\rho_n^{\mathrm e}(v)}
=
\frac{n^{-E(p)}}{n^{-E_0(p)}}
=
n^{E_0(p)-E(p)}
=
n^{G(v)}.
\]
Consequently, for every integer \(n\ge2\),
\[
\frac{c_v}{C^{\mathrm e}_v}n^{G(v)}
\le
\frac{r_n^*(v)}{r_n^{*,\mathrm e}(v)}
\le
\frac{C_v}{c^{\mathrm e}_v}n^{G(v)}.
\]
% lean: CausalSmith.Stat.FinitepHomogeneityDensegamma.preservation_ratio_bounds

\item If \(F(p)=1\), the exponent identity gives \(G(v)=0\). Substituting \(n^0=1\) into the quotient bounds gives
\[
\frac{c_v}{C^{\mathrm e}_v}
\le
\frac{r_n^*(v)}{r_n^{*,\mathrm e}(v)}
\le
\frac{C_v}{c^{\mathrm e}_v}
\qquad(n\ge2).
\]
Both bounding constants are strictly positive by \cref{obj:thm:heavy-tail-comparison,obj:thm:oracle-frontier-resolved}.
% lean: CausalSmith.Stat.FinitepHomogeneityDensegamma.exact_propensity_preservation_boundary

\item Suppose \(F(p)<1\). The exponent identity gives
\[
G(v)=
\frac{2\gamma q}{(2\gamma+q)D_p}\bigl(1-F(p)\bigr)>0.
\]
A positive power of the natural sample size tends to infinity, and multiplication by the positive lower multiplier preserves this limit:
\[
n^{G(v)}\longrightarrow+\infty,
\qquad
\frac{c_v}{C^{\mathrm e}_v}n^{G(v)}
\longrightarrow+\infty.
\]
For all sufficiently large natural numbers \(n\), we have \(n\ge2\), so the quotient lower bound gives
\[
\frac{c_v}{C^{\mathrm e}_v}n^{G(v)}
\le
\frac{r_n^*(v)}{r_n^{*,\mathrm e}(v)}.
\]
Comparison with the diverging lower bound proves
\[
\frac{r_n^*(v)}{r_n^{*,\mathrm e}(v)}
\longrightarrow+\infty.
\]
% lean: CausalSmith.Stat.FinitepHomogeneityDensegamma.preservation_ratio_tendsto

Moreover, channel selection gives \(E(p)=E_4(p)\), and hence
\[
G(v)=E_0(p)-E(p)=E_0(p)-E_4(p).
\]
Substitution into the quotient bounds proves
\[
\frac{c_v}{C^{\mathrm e}_v}n^{E_0(p)-E_4(p)}
\le
\frac{r_n^*(v)}{r_n^{*,\mathrm e}(v)}
\le
\frac{C_v}{c^{\mathrm e}_v}n^{E_0(p)-E_4(p)}
\qquad(n\ge2).
\]
% lean: CausalSmith.Stat.FinitepHomogeneityDensegamma.exact_propensity_preservation_boundary

\item We characterize the preservation region of
\cref{obj:def:preservation-region}. Suppose first that \(v\in\mathcal R\).
There is a real upper bound \(b_{\mathrm{sup}}\) satisfying
\[
\frac{r_n^*(v)}{r_n^{*,\mathrm e}(v)}
\le b_{\mathrm{sup}}
\qquad\text{for every integer }n\ge2.
\]
If \(F(p)<1\), the divergence just proved supplies an integer \(n\ge2\) such that
\[
b_{\mathrm{sup}}
<
\frac{r_n^*(v)}{r_n^{*,\mathrm e}(v)}
\le b_{\mathrm{sup}},
\]
a contradiction. Thus \(F(p)\ge1\).

Conversely, if \(v\in\mathcal V\) and \(F(p)\ge1\), then \(G(v)=0\), so
\[
\frac{r_n^*(v)}{r_n^{*,\mathrm e}(v)}
\le\frac{C_v}{c^{\mathrm e}_v}
\qquad\text{for every integer }n\ge2.
\]
This finite upper bound places \(v\) in \(\mathcal R\). Expanding \(F(p)\) therefore gives
\[
\mathcal R
=
\left\{v\in\mathcal V:
\frac{2\alpha}{p-1}
+\frac{\beta}{q}
+\frac{\alpha+\beta}{2\gamma}\ge1
\right\}.
\]
% lean: CausalSmith.Stat.FinitepHomogeneityDensegamma.exact_propensity_preservation_boundary

\item Finally, let \(\mathcal K\subseteq\mathcal V\) be an arbitrary compact set.
The compact-uniformity conclusion of
\cref{obj:thm:heavy-tail-comparison} supplies real constants
\(\underline c_{\mathcal K},\overline C_{\mathcal K}\) satisfying
\[
\underline c_{\mathcal K}>0,
\qquad
\underline c_{\mathcal K}\le c_v,
\qquad
C_v\le\overline C_{\mathcal K}
\qquad(v\in\mathcal K).
\]
The supplied-propensity upper multiplier is independent of \(v\) and positive. Define
\[
c_{\mathrm{pres},\mathcal K}
=
\frac{\underline c_{\mathcal K}}
{3\bigl(120+128\sqrt{160\sqrt2}\bigr)},
\qquad
C_{\mathrm{pres},\mathcal K}
=
64\sqrt3\,\overline C_{\mathcal K}.
\]
Then \(c_{\mathrm{pres},\mathcal K}>0\), and for every \(v\in\mathcal K\),
\[
c_{\mathrm{pres},\mathcal K}
\le\frac{c_v}{C^{\mathrm e}_v},
\qquad
64\sqrt3\,C_v
\le C_{\mathrm{pres},\mathcal K}.
\]
These are the required compact-uniform constants.
% lean: CausalSmith.Stat.FinitepHomogeneityDensegamma.exact_propensity_preservation_boundary
\end{enumerate}
\end{proof}
